\documentclass{article}

% if you need to pass options to natbib, use, e.g.:
% \PassOptionsToPackage{numbers, compress}{natbib}
% before loading nips_2017
%
% to avoid loading the natbib package, add option nonatbib:
% \usepackage[nonatbib]{nips_2017}

% to compile a camera-ready version, add the [final] option, e.g.:
\PassOptionsToPackage{sort&compress, numbers}{natbib}
%\usepackage{nips}
\usepackage[final]{nips}
\usepackage[utf8]{inputenc} % allow utf-8 input
\usepackage[normalem]{ulem}
\usepackage[T1]{fontenc}    % use 8-bit T1 fonts
\usepackage{hyperref}       % hyperlinks
\usepackage{url}            % simple URL typesetting
\usepackage{algorithm}
\usepackage{graphicx}
\usepackage{color}
\usepackage{algorithmic}
\usepackage{eqparbox}
\renewcommand\algorithmiccomment[1]{%
  \hfill\#\ \eqparbox{COMMENT}{#1}%
}
\usepackage{amsmath,amssymb}
\DeclareMathOperator*{\argmax}{argmax}
\usepackage{booktabs}       % professional-quality tables
\usepackage{amsfonts}       % blackboard math symbols
\usepackage{nicefrac}       % compact symbols for 1/2, etc.
\usepackage{microtype}      % microtypography
\usepackage{xcolor}
\usepackage{xspace}
\usepackage{verbatim}
\usepackage{graphicx}
\usepackage{subcaption}

\newcommand{\todo}[1]{\comments{\textcolor{red}{[TODO: #1]}}}
\newcommand{\done}[1]{\comments{\textcolor{green}{[DONE: #1]}}}
\newcommand{\add}[1]{\comments{\textcolor{red}{[TODO: insert #1]}}}
\newcommand{\jc}[1]{\comments{\textcolor{deepForestGreen}{[JC: #1]}}}
\newcommand{\blue}[1]{\comments{\textcolor{blue}{#1}}}
\newcommand{\vm}[1]{\comments{\textcolor{brown}{[VM: #1]}}}
\newcommand{\jl}[1]{\comments{\textcolor{magenta}{[Joel: #1]}}}
\newcommand{\ks}[1]{\comments{\textcolor{purple}{[Ken: #1]}}}
\newcommand{\fs}[1]{\comments{\textcolor{maroon}{[Felipe: #1]}}}
\newcommand{\rw}[1]{\comments{\textcolor{blue}{[Rui: #1]}}}
\definecolor{maroon}{rgb}{.5,0,0}
\definecolor{deepForestGreen}{rgb}{.2,.8,.04}
\definecolor{magenta}{rgb}{0.7,0,1}
\newcommand{\lowpriority}[1]{\comments{\textcolor{green}{[#1]}}}
\newcommand{\later}[1]{\comments{\textcolor{green}{[Later: #1]}}}
\newcommand{\maybe}[1]{\comments{\textcolor{lightgray}{[Maybe: #1]}}}
\newcommand{\removed}[1]{\comments{\textcolor{lightgray}{#1}}}

\newcommand{\comments}[1]{#1}

%\title{\todo{check title - under construction}The Paired Open-Ended Trailblazer: Continually Generating Diverse Challenges and Solutions in a Single Run}
\title{
Paired Open-Ended Trailblazer (POET): Endlessly Generating Increasingly Complex and Diverse Learning Environments and Their Solutions}


\author{
\textbf{Rui Wang} \hspace{5mm} \textbf{Joel Lehman} \hspace{5mm}  \textbf{Jeff Clune*}  \hspace{5mm}  \textbf{Kenneth O. Stanley*}\\
  Uber AI Labs\\
  San Francisco, CA 94103 \\
  \texttt{{ruiwang,joel.lehman,jeffclune,kstanley}@uber.com} \\
  *co-senior authors
}

\begin{document}

\maketitle
%\vspace{-0.8cm}
\begin{abstract}
While the history of machine learning so far largely
encompasses a series of problems posed by researchers and algorithms that learn their solutions, an important question is whether the problems themselves can be generated by the algorithm at the same time as they are being solved.  Such a process would in effect build its own diverse and expanding curricula, and the solutions to problems at various stages would become stepping stones towards solving even more challenging problems later in the process.  The \emph{Paired Open-Ended Trailblazer} (POET) algorithm introduced in this paper does just that: it pairs the generation of environmental challenges and the optimization of agents to solve those challenges.  It simultaneously explores many different paths through the space of possible problems and solutions and, critically, allows these stepping-stone solutions to transfer between problems if better, catalyzing innovation. The term \emph{open-ended} signifies the intriguing potential for algorithms like POET to continue to create novel and increasingly complex capabilities without bound.  We test POET in a 2-D bipedal-walking obstacle-course domain in which POET can modify the types of challenges and their difficulty. At the same time, a neural network controlling a biped walker is optimized for each environment.  The results show that POET produces a diverse range of sophisticated behaviors that solve a wide range of environmental challenges, many of which cannot be solved by direct optimization alone, or even through a direct-path curriculum-building control algorithm introduced to highlight the critical role of open-endedness in solving ambitious challenges.  The ability to transfer solutions from one environment to another proves essential to unlocking the full potential of the system as a whole, demonstrating the unpredictable nature of fortuitous stepping stones.  We hope that POET will inspire a new push towards open-ended discovery across many domains, where algorithms like POET can blaze a trail through their interesting possible manifestations and solutions.
%OLDER VERSION
%The goal of research on open-endedness is to create a single process that invents astronomical complexity for near-eternity, just like evolution on Earth. This branch of research has gone way beyond complexifying artificial lives in an artificial world, and has advanced with the introduction of diversity-promoting algorithms, such as novelty search (NS) and Quality diversity (QD), as well as coevolution-based frameworks. These advances already start to impact other fields within artificial intelligence (AI) that are also involving large, high-dimensional search spaces, including deep learning and reinforcement learning (RL). Here, we propose a novel open-ended coevolution algorithm, namely, \emph{Paired Open-Ended Trailblazer}, or POET, which is designed to facilitate an open-ended process of discovery within a single run. POET maintains one population of environments and one population of agents, and each environment is paired with an agent. While explicitly optimizing behavior of each agent within its paired environment, POET adopts a novel evolution strategies (ES) based transfer mechanism that enables an agent to periodically transfer its learned capability to other agents to boost the learning progress of that agent on its own paired environment. POET also adopts an ongoing divergent coevolutionary interaction between environments and agents to progressively create and add new environments to the population, and continually forges new paths to increasing skills and challenges. Through coevolving 2D walking courses and bipedal walking agents, we demonstrate the superior capability of POET on creating a wide range of diverse environments with increasing complexity as well as solvers with increasing capabilities and diverse behaviors at a level that a simple curriculum-based single lineage search process cannot reach. POET provides an interesting platform for RL algorithms to progressively generate diverse curriculum and learn interesting and diverse behaviors while mastering increasingly challenging tasks.
\end{abstract}
%\vspace{-0.5cm}



\section{Introduction} \label{intro} %\vspace{-0.3cm}
%\todo{Cite Giuseppe}\\
%\todo{Check 1,2}\\


Machine learning algorithms are often understood as tools for solving difficult problems.  Challenges like image classification and associated benchmarks like ImageNet \cite{Deng2009ImageNetAL} 
%\todo{cite Deng}\rw{done}
are chosen by people in the community as solvable and researchers focus on them to invent ever-more high-performing algorithms.  For example, ImageNet was proposed in 2009 %\todo{insert year}\rw{done} 
and modern deep neural networks such as \textit{ResNet} \cite{he:arxiv15} began to beat humans in 2015 \cite{russakovsky:arxiv14,he2:arxiv15}. 
%\ks{is this right?  can we cite something more definitive than an arxiv paper?}\rw{verified and more definitive citation are added}.  
In reinforcement learning (RL) \cite{sutton1998reinforcement}, learning to play Atari was proposed in 2013 \cite{bellemare2013arcade}
%\todo{insert year and cite Bellemare}\rw{done} 
and in recent years deep reinforcement learning (RL) algorithms reached and surpassed human play across a wide variety of Atari games \cite{mnih2015,espeholt2018impala,horgan2018distributed,hessel2018multi,kapturowski2018recurrent,goexplore_blog}.
%\todo{we need to cite more recent/best Atari papers, perhaps Impala, Apex, but even more impressive ones..whatever is the current SOTA on average across Atari games. Is it PopArt? Ask Felipe. Also cite the Go-Explore blog post here}\rw{done. Added IMPALA, Ape-X, Recurrent Ape-X, PopArt, Go-Explore}
The ancient game of Go has been a challenge in AI for decades, and has been the subject of much recent research: AlphaGo and its variants can now reliably beat the world's best human professional Go players \cite{SilverHuangEtAl16nature, silver2017mastering, silver2018general}.
%\todo{cite AlphaGo Zero}\rw{Done.}
Each such achievement represents the pinnacle of a long march of increasing performance and sophistication aimed at solving the problem at hand.

As compelling as this narrative may be, it is not the only conceivable path forward.  While the story so far relates a series of challenges that are conceived and conquered by the community algorithm by algorithm, in a more exotic alternative 
%\sout{not yet widely discussed} 
only rarely discussed (e.g.\ \cite{forestier2017intrinsically,schmidhuber2013powerplay}), the job of the
algorithm could be to conceive \emph{both} the challenges and the solutions at the same time.  Such a process offers the novel possibility that the march of progress, guided so far by a sequence of problems
conceived by humans, could lead \emph{itself} forward, pushing the boundaries of performance autonomously and indefinitely.  In effect, such an algorithm could continually invent new environments that pose novel problems just hard enough to challenge current capabilities, but not so hard that all gradient is lost.  The environments need not arrive in a strict sequence either; they can be invented in parallel and asynchronously, in an ever-expanding tree of diverse challenges and their solutions.  

We might want such an \emph{open-ended} \cite{standish:ijcia03,langdon:metas05,bedau:alife08, ken:oriley17,taylor:alife16} %\todo{we need more cites than just the oreilly article}\rw{added}\ks{moved some of these to the paragraph below about the "small community"; still need some more: Mark Bedau should be cited}\rw{added one more paper from Mark Bedau} 
process because the chain that leads from the capabilities of machine learning today to e.g.\ general human-level intelligence could stretch across vast and inconceivable paths of stepping stones.  There are so many directions we could go, and so many problems we could tackle, that the curriculum that leads from here to the farthest reaches of AI is beyond the scope of our present imagination.  Why then should we not explore algorithms that \emph{self-generate} their own such curricula?  If we could develop and refine such algorithms, we might find a new approach to progress that relies less upon our own intuitions about the right stepping stones and more upon the power of automation.

In fact, the only process ever actually to achieve intelligence at the human level, natural evolution, \emph{is} just such a self-contained and open-ended curriculum-generating process.  Both the problems of life, such as reaching and eating the leaves of trees for nutrition,
%\jl{Maybe reaching and eating the leaves -- otherwise reads a bit ambiguously at first.}, 
and the solutions, such as giraffes, are the products of the same open-ended process.  And this process unfolds not as a single linear progression,
%\jl{Maybe: this process unfolds not as a single linear progression},
but rather involves innumerable parallel and interacting branches radiating for more than a billion years (and is still going).  Nature is also a compelling inspiration because it has avoided convergence or stagnation, and continues to produce novel artifacts for beyond those billion years.  An intriguing question is whether it is possible to conceive an algorithm whose results would be worth waiting a billion years to see.  Open-ended algorithms in their most grandiose realization would offer this possibility.  

While a small community within the field of artificial life \cite{standish:ijcia03,langdon:metas05, lehman:alife08,graening:ppsn10,soros:alife14, soros2:alife14, brant:mcc17, ken:oriley17,ray:alife91,taylor:alife16}
%\todo{cite}\rw{added} 
has studied the prospects of open-ended computation for many years, in machine learning the prevailing assumption that algorithms should be \emph{directed} (e.g.\ towards the performance objective) has pervaded the development of algorithms and their evaluation for many years. Yet this explicitly directed paradigm has begun to soften in recent years. Researchers have begun to recognize that modern learning algorithms' hunger for data presents a long-term problem:  the data available for progress diminishes as the appropriate tasks for leveling up AI competencies increase in complexity. 
%\jl{long sentence, a bit hard to follow, especially at its culmination}.  
This recognition is reflected in systems based on \emph{self-play} (which is related to \emph{coevolution} \cite{ficici:alife98, wiegand:gecco01,popovici:hnc12}), such as competitive two-player RL competitions \cite{selfplay:arxiv18,silver2017mastering} wherein the task is a function of the competition, which is itself changing over time. Generative adversarial networks (GANs) \cite{goodfellow_generative_2014}, 
%\todo{cite}\rw{added}, 
in which networks interact as adversaries, similarly harness a flavor of self-play and coevolution \cite{moran:arxiv18} %\jl{Maybe cite the recent paper on connections between GANs and coevolution from Pollack? Can't remember if it would be appropriate here or not.} 
to generate both challenges and solutions to those challenges. 

%Original revision discussion to add new references about generating problems and solutions
%\rw{added: \ks{While, not explicitly motivated by open-endedness, recognition of the importance of self-generated curricula} \sout{Such recognition} is also reflected in recent advances in automatic curriculum learning for RL, where the intermediate goals of curricula are automatically generated and selected. \sout{, e.g., either by}\ks{A number of such approaches to automatic curriculum building have been proposed.  For example, the curriculum can be produced by a generator network optimized using adversarial training to produce tasks at the appropriate level of difficulty for the agent \cite{florensa2018automatic}, by applying noise in action space to generate a set of feasible start states increasingly far from the goal (i.e.\ a reverse curriculum) \cite{pmlr-v78-florensa17a}, or through intrinsically motivated goal exploration processes (IMGEPs) \cite{forestier2017intrinsically}.} The POWERPLAY algorithm offers an alternative approach to formulating continual, concurrent searches for new tasks through an increasingly general problem solver \cite{schmidhuber2013powerplay}. In Teacher-Student Curriculum Learning \cite{matiisen2017teacher},  the teacher automatically select for students sub-tasks with the highest slope of the learning curve of the student. Ideas from the field of \emph{procedural content generation} (PCG) \cite{togelius:ieeetciaig11,shaker2016procedural}, where methods are developed to generate diverse sets of levels and other contents for games, can also generate progressive curricula that help improve generality of deep RL agents \cite{justesen2017procedural}.}     

While not explicitly motivated by open-endedness, recognition of the importance of self-generated curricula is also reflected in recent advances in automatic curriculum learning for RL, where the intermediate goals of curricula are automatically generated and selected. A number of such approaches to automatic curriculum building have been proposed.  For example, the curriculum can be produced by a generator network optimized using adversarial training to produce tasks at the appropriate level of difficulty for the agent \cite{florensa2018automatic}, by applying noise in action space to generate a set of feasible start states increasingly far from the goal (i.e.\ a reverse curriculum) \cite{pmlr-v78-florensa17a}, or through intrinsically motivated goal exploration processes (IMGEPs) \cite{forestier2017intrinsically}. The POWERPLAY algorithm offers an alternative approach to formulating continual, concurrent searches for new tasks through an increasingly general problem solver \cite{schmidhuber2013powerplay}. In Teacher-Student Curriculum Learning \cite{matiisen2017teacher},  the teacher automatically selects sub-tasks that offer the highest slope of learning curve for the student. Ideas from the field of \emph{procedural content generation} (PCG) \cite{togelius:ieeetciaig11,shaker2016procedural}, where methods are developed to generate diverse sets of levels and other contents for games, can also generate progressive curricula that help improve generality of deep RL agents \cite{justesen2017procedural}.    

The ultimate implication of this shift towards less explicitly-directed processes is to tackle the full grand challenge of open-endedness: Can we ignite a process that on its own unboundedly produces increasingly diverse and complex challenges at the same time as solving them?  And, assuming computation is sufficient, can it last (in principle) forever?  

The \emph{paired open-ended trailblazer} (POET) algorithm introduced in this paper aims to confront open-endedness directly, by evolving a set of diverse and increasingly complex environmental challenges at the same time as collectively optimizing their solutions.  A key opportunity afforded by this approach is to attempt \emph{transfers} among different environments.  That is, the solution to one environment might be a stepping stone to a new level of performance in another, which reflects our uncertainty about the stepping stones that trace the ideal curriculum to any given skill.  By radiating many 
paths of increasing challenge simultaneously, a vast array of potential stepping stones emerges from which progress in any direction might originate.  POET owes its inspiration to recent algorithms such a minimal criterion coevolution (MCC) \cite{brant:mcc17}, which shows that environments and solutions can effectively co-evolve, novelty search with local competition \cite{lehman:gecco11}, MAP-Elites \cite{mouret2015illuminating,cully:nature15}, Innovation Engines \cite{nguyen2016innovative}, which exploit the opportunity to transfer high-quality solutions from one objective among many to another, and the 
combinatorial multi-objective evolutionary algorithm (CMOEA), which extends the Innovation Engine to combinatorial tasks \cite{huizinga2016does, Joost:arxiv18}.
POET in effect combines the ideas from all of these approaches to yield a new kind of open-ended algorithm that optimizes, complexifies, and diversifies as long as the environment space and available computation will allow.

In this initial
%\jl{redundant?}\jc{ I don't think so because it says that we don't think this is a definitive exhaustive introductory coverage of this idea, but instead suggests of this is a preliminary look into poet and we're going to do a lot more work even before we declare it THE intro} 
introduction of POET, it is evaluated in a simple 2-D bipedal walking obstacle-course domain in which the form of the terrain is evolvable, from a simple flat surface to heterogeneous environments of gaps, stumps and rough terrain.  The results establish that (1) solutions found by POET for challenging environments cannot be found directly on those same environmental challenges by optimizing on them only from scratch; 
(2) neither can they be found through a curriculum-based process aimed at gradually building up to the same challenges POET invented and solved;
(3) periodic transfer attempts of solutions from some environments to others--also known as ``goal switching'' \cite{nguyen2016innovative}--is important for POET's success; (4) a diversity of challenging environments are both invented and solved in the same single run.
%\ks{need to remember to insert a fourth one: that you also cannot rediscover the same solutions through a canned curriculum -- I haven't thought through the wording yet but don't want to forget that here}. 
While the 2-D domain in this paper offers an initial hint at the potential of the approach for achieving open-ended discovery in a single run, it will be exciting to observe its application to more ambitious problem spaces in the future.  

The paper begins with background on precedent for open-endedness and POET, turning next to the main algorithmic approach.  It finally introduces the 2-D obstacle-course domain and presents both qualitative an quantitative results.



\begin{comment}

%\ks{revised version above - OLD VERSION BELOW}

In recent years, there have been tremendous incredible achievements from machine learning and artificial intelligence (AI), especially through techniques that make use of large and deep artificial neural networks, widely recognized as \textit{deep learning} \cite{DBLP:journals/nature/LeCunBH15, Goodfellow-et-al-2016}. In numerous areas, deep neural networks have already demonstrated astonishing capabilities that are comparable to or outperform human performance. In image classification, modern deep neural networks such as \textit{ResNet} \cite{he:arxiv15} started to beat human record ever since back in 2015. In natural language processing (NLP), researchers recently presented a deep bidirectional Transformer model \cite{Devlin:arxiv18} that redefines the state of the art for 11 difficult tasks, even surpassing human performance in the challenging area of question answering \cite{D16-1264}. In the area of reinforcement learning (RL) \cite{sutton1998reinforcement}, deep RL algorithms reached and surpassed a level comparable to that of a professional human gametester across a wide variety of Atari games \cite{mnih2015}. More recently, AlphaGo and its variant \cite{SilverHuangEtAl16nature, silver2017mastering} invented at DeepMind beat the world's best human professional Go players. OpenAI's team of five neural networks, namely, \textit{OpenAI Five} \cite{OpenAI_dota}, has started to defeat human teams at Dota 2, a popular video game that requires sophisticated strategies and planning over very long time horizon and collaboration between players, which many believe are extremely difficult for deep RL agents to master. The list can go on and on. In fact, despite neural network techniques being notoriously temperamental, there is great confidence that this connectionist paradigm will continue to collect trophy victories. 

However, there is a strengthening voice in the AI community that has questioned the power of current techniques to lead the goal of human-level general intelligence across a general range of tasks. Though steady progress continues, no such general digital intelligence seems within our reach yet \cite{ken:oriley17}. Interestingly, an often understated fact is, there already exists a format of human-level general intelligence, that is, the intelligence of ourselves, and there is a process or an algorithm that has created it, that is, evolution, or more precisely, \textit{open-ended} evolution. Indeed, evolution on Earth is an open-ended process: the first life form, emerged billions of years ago, likely something resembling a single prokaryotic cell, over the next few billion years, radiated into the full diversity of life on Earth we appreciate today, including ourselves and our intelligence. Think about it computationally—evolution on Earth is like a single run of a single algorithm that invented all of our nature with astronomical complexity. What makes this algorithm amazingly different from almost all machine learning algorithms we are familiar with is that it never seems to end. Evolution can be viewed as a never-ending algorithm that tirelessly invents ever-greater complexity and novelty across incomprehensible spans of time, while human-level intelligence on Earth is one of its many products. In other words, open-endedness could be a prerequisite to AI, or at least a promising path to it. 

The rich interconnection of AI and open-endedness certainly goes beyond that open-endedness could be a force to discover intelligence. Open-endedness could also be a core component of intelligence itself. This is based on the observation that human intelligence is indeed open-ended: we don't stop at only optimizing our minds and skills to perform existing tasks. Instead we keep inventing new tasks and trying to identify new problems to solve. Besides, we are also playful and like to be creative simply to stimulate ourselves, even if no particular problem is solved, such as in art and music. If we believe human-level intelligence needs to have endless creativity like ourselves, open-endedness should be one of innovation engines of such intelligence.

Because of the inspiration from natural evolution, the study of open-ended algorithms has focused primarily on evolutionary algorithms (EAs). However, many regular EAs are not aimed to be open-ended at all. Converging to a solution or not, they all run for a while, and eventually exhaust the search space of anything new. Historically, many open-endness related research is done by a small community of scientists within the field called \textit{artificial life}, where they focused on creating an artificial world and hope to observe complexity explosion of diverse artificial creatures. (Readers can refer to \cite{ken:oriley17} for a comprehensive literature review.) Beyond that, the scope of research into open-endedness further broadened with the introduction of \textit{novelty search} (NS) \cite{lehman} and \textit{quality diversity} (QD) \cite{cully:nature15, mouret2015illuminating, pugh:frontiers16} algorithms, which will be explained in Sec. \ref{prelim}. NS and QD have addressed finding new solutions to existing challenges open-endedly, but they do not inherently generate new kinds of challenges to be solved. More recent open-ended systems \cite{brant:mcc17} generate both through coevolution \cite{popovici:hnc12}, which is what happens when different individuals in a population interact with each other while they are evolving—a field closely relevant to open-endedness. Finally, research advances mentioned above related to solving the challenge of open-endedness already start to profoundly impact other fields within AI that are also involving search within large, high-dimensional spaces, including deep learning. For example, NS and QD have been hybridized with evolution strategies (ES) \cite{es} to improve exploration on sparse or deceptive deep RL tasks \cite{conti:nses_nips2018}. Other recent high-profile results such as generative adversarial networks (GANs) \cite{goodfellow_generative_2014} and self-play learning in board games \cite{silver2017mastering} and competitions between reinforcement-learning robots \cite{selfplay:arxiv18} have strongly benefited from coevolutionary-like dynamics.
rongly benefited from coevolutionary-like dynamics.

Inspired by the ideas above, in this paper, we propose a novel open-ended coevolution framework, namely, \textit{Progressive Open-Ended Trailblazer}, or POET. POET starts with a trivial task and progressively adds new tasks and their solvers. Each task, when created, has a dedicated solver being optimized to master it. The solvers are optimized by a novel multi-objective ES with a novel ES-based transfer mechanism that the solver for one task can periodically transfer its learned capability to its peer to boost the learning progress of that peer on its own task. The reproduction of a task (i.e., to create a new task) is controlled by minimal criteria with respect to the current capabilities of solvers. Through coevolution, POET is capable of creating a wide range of tasks with increasing complexity as well as solvers with increasing capabilities and diverse behaviors that can solve those tasks. We demonstrate the superior capability of POET through coevolution of 2D bipedal walkers and walking courses. POET provides an interesting platform for RL algorithms to progressively generate diverse curriculum with increasingly challenging tasks and learn interesting and diverse behaviors.
%\vspace{-0.43cm}

\end{comment}

\section{Background} \label{prelim}
%\vspace{-0.3cm}

This section first reviews several approaches that inform the idea behind POET of tracking and storing stepping stones, which is common in evolutionary approaches to behavioral diversity. It then considers a recently proposed open-ended coevolutionary framework based on the concept of minimal criteria and, lastly, evolution strategies (ES) \cite{es}, %\todo{cites}\rw{done} 
which serves as the optimization engine behind POET in this paper (though in principle other reinforcement learning algorithms could be substituted in the future).  
\subsection{Behavioral Diversity and Stepping Stones}
\label{CMOEA overview}
%\vspace{-0.2cm}

A common problem of many stochastic optimization and search algorithms is becoming trapped on local optima, which prevents the search process from leaving sub-optimal points and reaching better ones. In the context of an open-ended search, becoming trapped is possible in more than one way: In one type of failure, the domains or problems evolving over the course of search could stop increasing their complexity and thereby stop becoming increasingly interesting.  Alternatively, the simultaneously-optimizing \emph{solutions} could be stuck at sub-optimal levels and fail to solve challenges that are solvable. Either scenario can cause the search to stagnate, undermining its open-endedness. 

%\ks{fix all quotation marks in document to be correct for latex – refer to latex manual for correct char sequences for left and right quotation marks}\rw{DONE: ``''}
Population-based algorithms going back to novelty search \cite{lehman} that encourage \emph{behavioral} (as opposed to genetic) diversity \cite{cully:nature15, mouret2015illuminating,pugh:frontiers16, huizinga2016does, nguyen2016innovative, Joost:arxiv18} have proven less susceptible to local optima, and thus naturally align more closely with the idea of open-endedness as they focus on \emph{divergence} instead of \emph{convergence}. These algorithms are based on the observation that the path to a more desirable or innovative solution often involves a series of waypoints, or \emph{stepping stones}, that may \emph{not} increasingly resemble the final solution, and are not known ahead of time. Therefore, divergent algorithms reward and preserve diverse behaviors to facilitate the preservation of potential stepping stones, which could pave the way to both a solution to a particular problem of interest and to genuinely open-ended search. 

%\ks{make sure all acronyms are defined first before using them}\rw{double-checked. But may need to revisit in case we rewrite Introduction}
In the canonical example of \emph{novelty search} (NS)
%\todo{acronym already defined?}\rw{done}
\cite{lehman}, individuals are selected purely based on how different their behaviors are compared to an archive of individuals from previous generations. In NS, the individuals in the archive determine the \textit{novelty} of a solution, reflecting the assumption that genuinely novel discoveries are often stepping stones to further novel discoveries.  Sometimes, a chain of novel stepping stones even reaches a solution to a problem, even though it was never an explicit objective of the search.  NS was first shown effective in learning to navigate deceptive mazes and biped locomotion problems  \cite{lehman}. 
%\todo{cite}\rw{done}

%\ks{note that periods come *before* end quotes in American English.  That is: That is a “cat.” Is correct.  That is a “cat”. Is incorrect.  British English does it the other way, but we’re in the US so can follow American conventions.}\rw{Done}
Other algorithms are designed to generate, retain, and utilize stepping stones more explicitly. To retain as much diversity as possible, quality diversity (QD) algorithms
%\ks{was the acronym QD ever defined?}\rw{yes, but may need to revisit if we rewrite the introduction a lot} algorithms 
\cite{lehman:gecco11,cully:nature15, mouret2015illuminating,pugh:frontiers16} 
%\ks{should also be citing Joel’s paper on novelty search with local competition, or NSLC, which is probably the original QD algorithm}\rw{Done} 
keep track of many different niches of solutions that are (unlike pure NS) being optimized 
%for different objectives or tasks, 
simultaneously 
%\jc{this makes it sound like the performance criteria are different, when usually it is behaviorally different ways of doing well according to one performance criterion}\ks{good catch, I reworded this sentence to be more accurate with respect to QD, but I think it was originally conflating CMOEA with all of QD in order to make the transirtion to the next paragraph fluid, which means it might still be a bit contorted -- hopefully it works now} 
%\ks{does introducing the new terminology “goals” here add useful information?  Or does it just compound the number of terms the reader needs to remember?}\rw{delete this terminology ``\textit{goals}''}, 
and in effect try to discover stepping stones by periodically testing the performance of offspring from one niche
%on tasks or objectives of other niches
in other niches, a process referred to 
%\ks{if it’s really “often” referred to this way, then there should be at least a citation}\rw{delete often, citation added} 
as \emph{goal switching} \cite{nguyen2016innovative}.

Based on this idea, an approach called the Innovation Engine \cite{nguyen2016innovative} is able to evolve a wide range of images in a single run that are recognized with high-confidence as different image classes by a high-performing deep neural network trained on ImageNet \cite{alexnet2012}. In this domain, the Innovation Engine maintains separate niches for images of each class,  including the image that performed the best in that class (i.e.\ produced the highest activation value according to the trained deep neural network classifier). It then generates variants (perturbations/offspring) of a current niche  champion and not only checks whether such offspring are higher-performing in that niche (class), but also checks whether they are higher-performing in any other niche. Interestingly, the evolutionary path to performing well in a particular class often passes through other (oftentimes seemingly unrelated) classes that ultimately serve as stepping stones to recognizable objects. In fact, it is \emph{necessary} for the search to pass through these intermediate, seemingly unrelated stages to ultimately satisfy other classes later in the run \cite{nguyen2016innovative}.  More recently, to evolve multimodal robot behaviors, the \textit{Combinatorial Multi-Objective Evolutionary Algorithm} (CMOEA) \cite{huizinga2016does, Joost:arxiv18} builds on the Innovation Engine, but defines its niches based on different \emph{combinations of subtasks}.  This idea helps it to solve both a multimodal robotics problem with six subtasks as well as a maze navigation problem with a hundred subtasks, highlighting the advantage of maintaining a diversity of pathways and stepping stones through the search space when a metric for rewarding the most promising path is not known.  POET will similarly harness goal-switching within divergent search.

%\ks{The background is so far pretty heavy on evolutionary inspiration, which makes sense given its role in the new algorithm POET, but POET can hybridize with RL too, and we also may want to speak more explicitly to the broader RL community, in which case it could be useful to add some text on how the high-level ideas here on maintaining multiple pathways and stepping stones are actually somewhat generic and relevant even outside evolution, though they just happened to grow up within the field of evolutionary computation.}\rw{Propose to add discussion about two recent deep rl papers:} 
While the general ideas of promoting diversity and preserving stepping stones grew up initially within the field of evolutionary computation, they are nevertheless generic and applicable in fields such as deep reinforcement learning (RL). For example, a recent work utilizes diversity maximization to acquire complex skills with minimal supervision, which improves learning efficiency when there is no reward function. The results then become the foundation for imitation
learning and hierarchical RL \cite{eysenbach:arxiv18}. 
%\todo{something is wrong with the eysenbach:arxiv18 reference.  it lacks commas between authors and looks like its first author is ibarz and not eysenbach -- check carefully (and check all references for errors)}\rw{corrected. checking other references}
%and as a stepping stone for imitation learning and hierarchical RL. 
To address the problem of reward sparsity, \citet{savinov:arxiv18} extend the archive-based mechanism from NS that determines novelty in behaviors to determine novelty in the observation space and then to formulate the novelty bonus as an auxiliary reward. Most recently, the Go-Explore algorithm by Ecoffet et al. \cite{goexplore_blog} %\todo{cite blog post}\rw{done} 
represents a new twist on QD algorithms and is designed to handle extremely sparse and/or deceptive ``hard exploration'' problems in RL. It produced dramatic improvements over the previous state of the art on the notoriously challenging hard-exploration domains of Montezuma's Revenge and Pitfall. Go-Explore maintains an archive of interestingly different stepping stones (e.g. states of a game) that have been reached so far and how to reach them. It then returns to and explores from these stepping stones to find new stepping stones, and repeats the process until a sufficiently high-quality solution is found. 


\subsection{Open-Ended Search via Minimal Criterion Coevolution (MCC)}
\label{mcc overview}
%\vspace{-0.2cm}
Despite the recent progress of diversity-promoting algorithms such as NS and QD, they remain far from matching a genuinely open-ended search in the spirit of nature. 
%A primary problem is that these algorithms generally require a \emph{behavior characterization} (BC), which is a descriptor of individual behaviors in the search space.  The BC is what allows the measurement of novelty or the identification of niches.  Without it, there would be no opportunity to reward divergence.  Another unnatural requirement in algorithms like NS is the archive of points already visited.  Interestingly, the open-ended divergence of nature required neither a BC nor an archive of any type, suggesting that open-endedness might be encouraged or even explained by even more fundamental forces than explicit measures or records of behavior.
%\ks{Commented out the part criticizing needing a BC because we are now doing something
%like that with our diversity pressure for environments.  Trying a different angle here:}\rw{agree}
One challenge in such algorithms is that although they provide pressure for ongoing divergence
in the solution space, the \emph{environment} itself remains static, limiting the scope
of what can be found in the long run.
%\ks{Might make logical sense to say something brief about coevolution (and self-play) here, before starting on MCC. Coevolution is also inspired by open-endedness and also does not use a BC.  It explains where MCC came from and also shows why MCC is interesting as a new way of conceiving coevolution in a more QD-like setup.}\rw{Propose to add the following about coevolution:} 
In this spirit, one fundamental force for driving open-endedness could come from \emph{coevolution} \cite{popovici:hnc12}, which means that different individuals in a population interact with each other while they are evolving. Such multi-agent interactions have recently proven important and beneficial for systems with coevolutionary-like dynamics such as generative adversarial networks (GANs) \cite{goodfellow_generative_2014}, self-play learning in board games \cite{silver2017mastering}, and competitions between reinforcement-learning robots \cite{selfplay:arxiv18}.  Because the environment in such systems now contains other changing agents, the challenges are no longer static and evolve in entanglement with the solutions. 
%\jc{However, in most, if not all, coevolutionary systems, the abiotic (non-opponent-based) aspect of the environment remains fixed. No amount of co-evolution with a fixed task (e.g. pushing the other enemy off a pedestal, or beating them at Go) will produce general artificial intelligence that can write poetry, engage in philosophy, and invent math. An insight behind POET is that the environments themselves (including ultimately the reward function) must too change to drive truly open-ended dynamics. While much work in artificial life and evolutionary computation tries to create coevolutionary dynamics that will promote niche formation and the evolution of specialists that focus on such niches, an additional insight behind POET is to explicitly and separately evolve a population of environments to catalyze the production of diverse, complex niches. This spirit is motivated by a heuristic, which is that if we want an algorithm to generate some phenomenon, we should explicitly search for it instead of setting the stage and hoping that phenomenon emerges.} \ks{Stuck some of this useful additional argument into the start of the next paragraph; however, did not keep it exactly because it implies that this idea originates with POET when it actually originates (at least to a large extent) with MCC.  I want to be careful to give my student (Jonathan Brant) credit for what he did here.}

However, in most, if not all, coevolutionary systems, the abiotic (non-opponent-based) aspect of the environment remains fixed. No amount of co-evolution with a fixed task (e.g. pushing the other enemy off a pedestal, or beating them at Go) will produce general artificial intelligence that can write poetry, engage in philosophy, and invent math. 
Thus an important insight is that environments themselves (including ultimately the reward function) must too change to drive truly open-ended dynamics. 
Following this principle, an important predecessor to the POET algorithm in this paper is the recent \emph{minimal criterion coevolution} (MCC) 
%\ks{acronym defined?}\rw{Added.} 
algorithm \cite{brant:mcc17}, which explores an alternative paradigm for open-endedness through coevolution. In particular, it implements a novel coevolutionary framework that pairs a population of evolving \emph{problems} (i.e.\ environmental challenges) with a co-evolving population of \emph {solutions}.  
That enables new problems to evolve from existing ones and for new kinds of solutions to branch from existing solutions.  As problems and solutions proliferate and diverge in this ongoing coupled interaction, the potential for open-endedness arises.
%even without any explicit pressure for diversity.
%\jc{a dangerous comment because we sadly did add exactly this sort of pressure}

Unlike conventional coevolutionary algorithms that are usually divided between competitive coevolution, where individuals are rewarded according to their performance against other individuals in a competition \cite{ficici:alife98}, and cooperative coevolution, where instead fitness is a measure of how well an individual performs in a cooperative team with other evolving individuals \cite{wiegand:gecco01}, MCC introduce a new kind of coevolution that evolves two interlocking populations whose members earn the right to reproduce by satisfying a \emph{minimal criterion} \cite{soros:alife14, soros2:alife14, lehman:gecco10a} with respect to the other population, as both populations are gradually shifting simultaneously.  For example, in an example in \citet{ brant:mcc17} with mazes (the problems) and maze solvers (the solutions, represented by neural networks), the minimal criterion for the solvers is that they must solve at least one of the mazes in the maze population, and the minimal criterion for the mazes is to be solved by at least one solver.

The result is that the mazes increase in complexity and the neural networks continually evolve to solve them. In a single run without any behavior characterization or archive (which are usually needed in QD algorithms), MCC can continually generate novel, diverse and increasingly complex mazes and maze solvers as new mazes provide new opportunities for innovation to the maze solvers (and vice versa), hinting at open-endedness.  However, in MCC there is no force for \emph{optimization} within each environment (or maze in the example experiment).  That is, once a maze is solved there is no pressure to \emph{improve} its solution; instead, it simply becomes a potential stepping stone to a solution to another maze.  As a result, while niches and solutions proliferate, there is no pressure for \emph{mastery} within each niche, so we do not see the best that is possible. Additionally, in MCC problems have to be solvable by the \emph{current} population, so complexity can only arise through drift. As an example, imagine an environmental challenge $E'$ that is similar to a currently solved problem $E$, but a bit harder such that none of the current agents in the population can solve it. Further imagine that if we took some of the current agents (e.g. the agent that solves $E$) and optimized them to solve $E'$, at least one agent could learn to solve $E'$ in a reasonable amount of time. MCC would still reject and delete $E'$ if it was generated because none of the current population solves it right away. In contrast, in POET, we accept problems if (according to some heuristics) they seem likely to be improved upon and/or solved after some amount of dedicated optimization effort (so POET will likely keep $E'$), enabling a more direct and likely swifter path to increasingly complex, yet solvable challenges.

\subsection{Evolution Strategies (ES)} \label{evolution strategies details}%\vspace{-0.1cm}
In the POET implementation in this paper, ES plays the role of the optimizer (although other optimization algorithms should also work).  Inspired by natural evolution, ES \cite{rechenberg1978evolutionsstrategien} represents a broad class of population-based optimization algorithms. The method referred to here and subsequently as ``ES'' is a version of ES popularized by \citet{es} that was recently applied with large-scale deep learning architectures to modern RL benchmark problems. This version of ES draws inspiration from Section 6 of \citet{williams1992simple} (i.e.\ REINFORCE with multiparameter distributions), as well as from subsequent population-based optimization methods including Natural Evolution Strategies (NES) \cite{wierstra2008natural} and Parameter-Exploring Policy Gradients (PEPG) \cite{sehnke2010parameter}. More recent investigations have revealed the relationship of ES to finite difference gradient approximation \cite{lehman:gecco18esfd} and stochastic gradient descent \cite{zhang:arxiv17}. 

In the typical context of RL, we have an environment, denoted as $E(\cdot)$, and an agent under a parameterized policy whose parameter vector is denoted as $w$. The agent tries to maximize its reward, denoted as $E(w)$, as it interacts with the environment. In ES, $E(w)$ represents the stochastic reward experienced over a full episode of an agent interacting with the environment. Instead of directly optimizing $w$ to maximize $E(w)$, ES seeks to maximize the expected fitness over a population of $w$, $J(\theta) = \mathbb{E}_{w \sim p_{\theta}(w)}[E(w)]$, where $w$ is sampled from a probability distribution $p_{\theta}(w)$ parameterized by $\theta$. 
%and thus ES optimizes $\theta$ with stochastic gradient ascent\jc{this last clause is a non-sequitur}. 
Using  the log-likelihood trick \cite{williams1992simple} and estimating the expectation above over $n$ samples allows us to write the gradient of $J(\theta)$ with respect to $\theta$ as:
%\vspace{-0.15cm}
\begin{gather*}\label{eq:1}
\scalebox{1.}{$
\nabla_{\theta}J(\theta) \approx \frac{1}{n}\sum_{i=1}^{n}E(\theta_{i})\nabla_{\theta} \log p_{\theta}(\theta_{i}).
$}
\end{gather*}

Intuitively, this equation says that for every sample agent $\theta_i$, the higher the performance of that agent $E(\theta_i)$, the more we should move $\theta$ in the direction of the gradient that increases the likelihood of sampling that agent, which is $\nabla_{\theta} \log p_{\theta}(\theta_{i})$. 
%\jc{better?}
Following the convention in \citet{es}, $\theta_{i}, i = 1, \ldots, n,$ are $n$ samples of parameter vectors drawn from an isotropic multivariate Gaussian with mean $\theta$ and a covariance $\sigma^{2},$ namely, $\mathcal{N}(\theta,\,\sigma^{2}I)$. (The covariance can be fixed or adjusted over time depending on the implementation.) Note that $\theta_{i}$ can be equivalently obtained by applying additive Gaussian noise $\epsilon_{i} \sim \mathcal{N}(0,\,I)$ to a given parameter vector $\theta$ as $\theta_{i} = \theta + \sigma\epsilon_{i}$. With that notation, the above estimation of the gradient of $J(\theta)$ with respect to $\theta$ can be written as:
%\vspace{-0.1cm}
\begin{gather*}\label{eq:2}
\scalebox{1.}{$
\nabla_{\theta}J(\theta) \approx \frac{1}{n\sigma}\sum_{i=1}^{n}E(\theta + \sigma\epsilon_{i})\epsilon_{i}.
$}
\end{gather*}

%\jc{Why are the above two equations approximate instead of exact? Should we explain that?}\rw{because they are the estimated expectation over n samples. please see if the words before the first equation could serve as an explanation.}
Algorithm \ref{alg:es_step} illustrates the calculation of one optimization step of ES, which can efficiently be parallelized over distributed workers \cite{es}.
%\rw{propose to add to address Jeff's comment inside Algo \ref{alg:es_step}: Once calculated, the returned optimization step is added to the current parameter vector to obtain the next parameter vector.} 
Once calculated, the returned optimization step is added to the current parameter vector to obtain the next parameter vector. Note that implementations of ES including our own usually follow the approach of \citet{es} and rank-normalize $E(\theta + \sigma\epsilon_{i})$ before taking the weighted sum, which is a variance reduction technique.
%\todo{Rui: is this accurate? If so, please remove the jc tag and this todo}\rw{yes} 
Overall, ES has exhibited performance on par with some of the traditional, simple gradient-based RL algorithms on difficult RL domains (e.g. DQN \cite{mnih2015} and A3C \cite{a3c}), 
%\todo{Rui: is this accurate? If so, please remove the jc tag and this todo}\rw{yes, reference added}
including Atari environments and simulated robot locomotion \cite{es, conti:nses_nips2018}.
More recently, NS and QD algorithms have been shown possible to hybridize with ES to further improve its performance on sparse or deceptive deep RL tasks, while retaining scalability \cite{conti:nses_nips2018}, 
%\ks{update this citation to be the citation of the paper that is now in 2018 NeurIPS, which should subsume the arxiv version}\rw{done}\ks{This citation is still a bit broken.  First, the token conti:arxiv17 should not point to a NeurIPS paper -- it still should point to the arxiv version.  We should make a new token to point to the NeurIPS version and use that one.  Also, the reference itself is not correct for NeurIPS.  For example, it does not even show the name "NeurIPS" anywhere, and it's missing other info.  Please obtain proper bibtex from http://papers.nips.cc/paper/7750-improving-exploration-in-evolution-strategies-for-deep-reinforcement-learning-via-a-population-of-novelty-seeking-agents}\rw{Done. i pointed conti:arxiv17 back to the arxiv paper, added and used a new token conti:nses\_nips2018 based on the bibTex from the link, but delete fields (editors, publisher, url) to make th format consistent with other NeurIPS citation} 
providing inspiration for its novel hybridization within an CMOEA- and MCC-like algorithm in this paper.
%\ks{Should we put (NeurIPS) after "Advances in Neural Information Processing Systems" in the bibtex for NeurIPS papers?}\rw{i added NeurIPS.}

\begin{algorithm}[htp]
\caption{ES\_STEP}
\label{alg:es_step}
%\algsetup{linenosize=\tiny}
%\scriptsize
\begin{algorithmic}[1]
   \STATE {\bfseries Input:} an agent denoted by its policy parameter vector $\theta$, an environment $E(\cdot)$,  learning rate $\alpha$, noise standard deviation $\sigma$
  
   \STATE Sample $\epsilon_{1}, \epsilon_{2}, \ldots, \epsilon_{n} \sim \mathcal{N}(0, I)$
   \STATE Compute $E_{i} = E(\theta + \sigma\epsilon_{i})$ for $i = 1, \ldots, n$
   \STATE {\bfseries Return:} $\alpha \frac{1}{n\sigma} \sum_{i=1}^{n} E_{i} \epsilon_{i}$%\jc{we don't say what we do with this returned vector}\rw{add a sentence above to address this. please check.}

   
\end{algorithmic}
\end{algorithm}


%\vspace{-0.3cm}

%\vspace{-0.4cm}


\section{The Paired Open-Ended Trailblazer (POET) Algorithm} \label{algos} %\vspace{-0.3cm}
POET is designed to facilitate an open-ended process of discovery within a single run. It maintains a population of environments (for example, various obstacle courses) and a population 
%\jc{we also keep an archive of past environments and elites, right? we should mention that here too. or just switch to calling it all an archive and later talking about the difference between the active population and archive}\ks{what is this archive of past stuff used for? maybe that's what should be made clear} 
of agents (for example, neural networks that control a robot to solve those courses), and each environment is paired with an agent to form an \textit{environment-agent pair}. POET in effect implements an ongoing divergent coevolutionary interaction 
%\jc{makes it seem like the divergent co-evolutionary dynamic is within an environment-agent pair...need to clarify}
among all its agents and environments in the spirit of MCC \cite{brant:mcc17}, but with the added goal of explicitly optimizing the behavior of each agent within its paired environment in the spirit of CMOEA \cite{huizinga2016does, Joost:arxiv18}. 
It also elaborates on the minimal criterion in MCC by aiming to maintain only those newly-generated environments that are not too hard \emph{and} not too easy for the current population of agents.
%\jc{It also adds an entirely new feature in attempting, via heuristics, to estimate which newly generated environments deserve to be kept because they are likely to be not too easy and not too hard for the current population of solutions.}\ks{this is not entirely new -- MCC checks for being too hard, and could easily have also had an aspect of its MC to prevent things that are too easy.  It's more like a refinement on MCC.}
The result is a \emph{trailblazer} algorithm, one that continually forges new paths to both increasing challenges and skills
%, both generalized and specialized,
%\jc{I know what you mean here, but we are injecting this very complicated concept/observation without any explanation, so I do not think naive readers will have any idea what we mean by this nor understand its importance. We likely need to make this at least the subject of a separate sentence that follows, if not a separate paragraph.}
within a single run.  The new challenges are embodied by the new environments that are continually created, and the increasing skills are embodied by the neural network controllers attempting to solve each environment. Existing skills are harnessed both by optimizing agents paired with environments and by attempting to transfer current agent behaviors to new environments to identify promising stepping stones. 

The fundamental algorithm of POET is simple:  The idea is to maintain a list of active environment-agent pairs \texttt{EA\_List} that begins with a single starting pair  $(E^{\textrm{init}}(\cdot),\theta^{\textrm{init}})$, where $E^{\textrm{init}}$ is a simple environment (e.g.\ an obstacle course of entirely flat ground) and $\theta^{\textrm{init}}$ is a randomly initialized weight vector (e.g.\ for a neural network).  POET then has three main tasks that it performs at each iteration of its main loop: 

\begin{enumerate}
    \item generating new environments $E(\cdot)$ from those currently active, 
    \item optimizing paired agents within their respective environments, and 
    \item attempting to transfer current agents $\theta$ from one environment to another.  
\end{enumerate}

Generating new environments is how POET continues to produce new challenges.  To generate a new environment, POET simply mutates (i.e.\ randomly perturbs) the encoding (i.e.\ the parameter vector) of an active environment.  However, while it is easy to generate perturbations of existing environments, the delicate part is to ensure both that (1) the paired agents in the originating (parent) environments have exhibited sufficient progress to suggest that reproducing their respective environments would not be a waste of effort, and (2) when new environments are generated, they are not added to the current population of environments unless they are neither too hard nor too easy for the current population. %\jc{I think it would be much clearer to remove the rest of this sentence. \sout{for the agent paired with their parent environment (a clone of which will be the initial paired agent for that new environment)}}.  
Furthermore, priority is given to those candidate environments that are most novel, which produces a force for diversification that encourages many different kinds of problems to be solved in a single run.

These checks together ensure that the curriculum that emerges from adding new environments is smooth and calibrated to the learning agents.  In this way, when new environments do make it into the active population, they are genuinely stepping stones for continued progress and divergence.  The population of active environments is capped at a maximum size, and when the size of the population exceeds that threshold, the oldest environments are removed to make room (as in a queue).  That way, environments do not disappear until absolutely necessary, giving their paired agents time to optimize and allowing skills learned in them to transfer to other environments.

POET optimizes its paired agents at each iteration of the main loop.   The idea is that every agent in POET should be continually improving within its paired environment.  In the experiments in this paper, each such iteration is a step of ES, but any reinforcement learning algorithm could conceivably apply.  The objective in the optimization step is simply to maximize whatever performance measure applies to the environment (e.g.\ to walk as far as possible through an obstacle course).  The fact that each agent-environment pair is being optimized independently affords easy parallelization, wherein all the optimization steps can in principle be executed at the same time.

Finally, attempting \emph{transfer} is the ingredient that facilitates serendipitous cross-pollination: it is always possible that progress in one environment could end up helping in another.  For example, if the paired agent $\theta^A$ in environment $E^A(\cdot)$ is stuck in a local optimum, one remedy could be a transfer from the paired agent $\theta^B$ in environment $E^B(\cdot)$.  If the skills learned in the latter environment apply, it could revolutionize the behavior in the former, reflecting the fact that the most promising stepping stone to the best possible outcome may not be the current top performer in that environment \cite{stanley:book15, nguyen2016innovative,mouret2015illuminating}.  Therefore, POET continually attempts transfers among the active environments.  These transfer attempts are also easily parallelized because they too can be attempted independently.

More formally, Algorithm \ref{alg:poet} describes this complete main loop of POET. Each component is shown in pseudocode: creating new environments, optimizing each agent, and attempting transfers.  The specific implementation details for the subroutines MUTATE\_ENVS (where the new environments are created) and EVALUATE\_CANDIDATES (where transfers are attempted) are given in Supplemental Information. 

%(1) POET begins with an initial, often simple, environment $E^{\textrm{init}}(\cdot)$ paired a randomly initialized agent $\theta^{\textrm{init}}$. (2) POET periodically mutates the existing population of environments in the hope of creating and admitting new environments, and deletes old environments when necessary as described in Algorithm \ref{alg:mutate_niches}. (3) Each agent is independently optimized via ES to maximize its reward in its paired environment, while periodically POET attempts mutual transfer among all agents and environments.\jc{ it's not quite clear when we are asking the reader to stop reading the main text and go read the pseudocode. also it's not clear how these three phases map onto the pseudocode. Could we add comments to the code to demarcate where these three phases are in relation to the code?}

\begin{algorithm}[htp]
\caption{POET Main Loop}
\label{alg:poet}
%\algsetup{linenosize=\tiny}
%\scriptsize
%\ks{comments might be helpful in this code.  for example, each for loop could have a brief comment on what it's supposed to be doing.}\rw{added}
\begin{algorithmic}[1]
   \STATE {\bfseries Input:} initial environment $E^{\textrm{init}}(\cdot)$, its paired agent denoted by policy parameter vector $\theta^{\textrm{init}}$, learning rate $\alpha$, noise standard deviation $\sigma$, iterations $T$, mutation interval $N^{\textrm{mutate}}$, transfer interval $N^{\textrm{transfer}}$
   \STATE {\bfseries Initialize:} Set \texttt{EA\_list} empty 
   \STATE Add $(E^{\textrm{init}}(\cdot)$, $\theta^{\textrm{init}})$ to \texttt{EA\_list}
   \FOR{$t=0$ {\bfseries to} $T-1$}
   \IF{$t > 0$ {\bfseries and} $t \mod N^{\textrm{mutate}} = 0$}
   \STATE $\texttt{EA\_list} =$ MUTATE\_ENVS$(\texttt{EA\_list})$ \COMMENT{\bfseries new environments created by mutation}
   \ENDIF
   \STATE $M = \textrm{len}(\texttt{EA\_list})$
   \FOR{$m=1$ {\bfseries to} $M$} 
   \STATE $E^{m}(\cdot)$, $\theta^{m}_{t}$ = \texttt{EA\_list}[m]
   \STATE $\theta^{m}_{t+1} = \theta^{m}_{t} +$ ES\_STEP$(\theta^{m}_{t}, E^{m}(\cdot), \alpha, \sigma)$ \COMMENT{\bfseries each agent independently optimized} 
   \ENDFOR
   \FOR{$m=1$ {\bfseries to} $M$} 
   \IF{$M > 1$ {\bfseries and} $t \mod N^{\textrm{transfer}} = 0$} 
   \STATE $\theta^{\textrm{top}} = $ EVALUATE\_CANDIDATES$(\theta^{1}_{t+1}$, \ldots, $\theta^{m-1}_{t+1}$, $\theta^{m+1}_{t+1}$, \ldots, $\theta^{M}_{t+1}$, $E^{m}(\cdot)$, $\alpha$, $\sigma)$
   \IF{$E^{m}(\theta^{\textrm{top}}) > E^{m}(\theta^{m}_{t+1})$}
   \STATE $\theta^{m}_{t+1} = \theta^{\textrm{top}}$ \COMMENT{\bfseries transfer attempts}
   \ENDIF
   \ENDIF
   \STATE $\texttt{EA\_list}[m] = (E^{m}(\cdot), \theta^{m}_{t+1})$  
   
   \ENDFOR
   \ENDFOR
   


   
\end{algorithmic}
\end{algorithm}

%\ks{Added this paragraph to try to end the technical description of the algorithm with something more conceptual and motivational before switching to the experiment.}
%A unique aspect of POET is that new environments are continually generated.  Those that are not discarded are sufficiently close to their predecessors that they are not too difficult, but also not too trivial.  At the same time, an optimization algorithm (in this case ES acting as a reinforcement learner) improves the performance of the current policies within each environment niche.  What makes this optimization process unique is that policies from individual niches are continually offering their skills in other niches, to see if they might offer even better options\jc{ that's not entirely correct both Innovation engines and cmoea do this too, as do most QD algorithms like map-elites and NSLC.  point out those connections here. what is unique is that we're pairing that goal switching dynamic with the generation of new niches for the first time ever}.  This kind of cross-environment transfer respects the idea that the most promising stepping stone to the best possible outcome may not be the current top performer in that environment \cite{stanley:book15, nguyen2016innovative,mouret2015illuminating}.  As results in the experiment introduced next will show, it turns out that such transfer is essential to achieving the best performance in many environments.  Furthermore, as a result of its coevolutionary dynamic, POET generates a wide variety of increasingly challenging environments \emph{and} solutions to all of them, all in a single run, hinting at its potential for open-endedness.

% \ks{Missing in this section: We should very briefly mention the parallel computational platform.  The details are not important, but it is important to let the reader know we built a platform that full explouts parallelism, and that such parallelism is essential for this kind of divergent system to be computationally feasible.}\rw{agree. added below}

As noted above, the independence of many of the operations in POET, such as optimizing individual agents within their paired environments and attempting transfers, makes it feasible to harness the power of many processors in parallel.  In the implementation of the experiment reported here, each run harnessed 256 parallel CPU Cores.
Our software implementation of POET, which will be released as open source code shortly,
%TODO: ADD link when available!
allows seamless parallelization over any number of cores: workers are managed through Ipyparallel to automatically distribute all current tasks so that requests for specific operations and learning steps can be made without concern for how they will be distributed.

Provided that a space of possible environmental challenges can be encoded, the hope is that the POET algorithm can then start simply and push outward in parallel along an increasingly challenging frontier of challenges, some of which benefit from the solutions to others.  The next section describes a platform designed to test this capability.

%For a divergent system like POET to be computationally feasible, exploiting parallelism in the implementation is essential. POET is implemented to fully exploit the parallelism inside its algorithm, and run in this work in practice on up to 256 workers\jc{ the grammar in this sentence is not right but I'm not sure exactly what it is trying to say so I do not know how to fix it}\jl{I get the gist of what you're saying but it is very hard to parse as written -- and maybe should say more about what a worker is -- a worker machine, a worker CPU core?}.  The mutation and evaluation of environments can be easily distributed among workers. ES-based agent optimization is highly parallelizable \cite{es}, and the computations of periodic transfer experiments can also be executed as ES optimization steps 
%\ks{something still seems wrong here -- mutations of environments do not happen through ES, so how can "all be executed as one or more ES optimization steps"?}\rw{agree. corrected.}, 
%and therefore can be parallelized in a similar manner.\jc{ I think this paragraph should be Rewritten to State exactly what parallelized. there's talk of what could be done and how ES is parralelizable  in general without telling us whether or not we do that in this paper and somehow the whole is less than the sum of the parts in this paragraph. I end up more confused and only vaguely know that some parallelization is occurring but not exactly what and how}\jl{Second that this paragraph needs some remolding.}

%POET maintains a population of environments, each of which is paired with an agent optimized by ES. 
%by separating workers into \emph{rollout workers}, which are dedicated to running batches for any given agents on batches of any given environments.  The workers then return the cumulative rewards, and optimizer workers, which perform aggregating of returns and weight updates. 
%\ks{the rest of this sentence is incoherent: wheat does "and optimizer works, which..." mean?  needs rewording} \ks{I have a feeling this paragraph could be shorter.  It has too many unnecessary details.  the main point is just that it's all parallelized.  We don't really need separate explanations of each kind of step or scenario.}\rw{agree, commented out some unnecessary details and rewritten.}

%The first component of POET that is introduced is a novel ES-based transfer mechanism based on Algorithm \ref{alg:es_step}, which is simple, but effective. Second, we describe how environments are generated by mutating current environments in the population and checking their potential for supporting progress
%\jc{if we switch to archive instead of population, as I recommend, then we need to search for the word population and update the language throughout}, which forms the basis of the coevolutionary framework.
%\jc{I don't love this last clause. The framework is based on many interacting pieces, not just these two. For example, which environments we accept.}. 
%Finally, the overall POET framework is unified to conclude the section. 
%\jl{I agree with Jeff that these sections are a slog -- could we begin with the conceptual unfolding of POET rather than conclude with it? I would like to be able to skim the methods section and somewhere get a high-level sense of how it works without diving into the nitty gritty.}
%\jc{this is an odd, confusing paragraph. I get two (of many) components described, then another overview?}   


%\vspace{-0.4cm}
\section{Experiment Setup And Results} \label{experiments} %\vspace{-0.3cm}

%HERE

%\ks{Added this paragraph to motivate the experiment.  We should explain why we chose this setup of all possible domains we could choose.}
An effective test of POET should address the hypothesis that it can yield an increasingly challenging set of environments, many with a satisfying solution, all in a single run.  Furthermore, we hope to see evidence for the benefit of cross-environment transfers.  The experimental domain in this work is a modified version of the ``Bipedal Walker Hardcore'' environment of the OpenAI Gym \cite{brockman}.  Its simplicity as a 2-D walking domain with various kinds of possible terrain makes it easy to observe and understand qualitatively different ambulation strategies  simply by viewing them.  Furthermore, the environments are easily modified, enabling numerous diverse obstacle courses to emerge to showcase the possibilities for adaptive specialization and generalization. Finally, it is relatively fast to simulate, allowing many long experiments compared to the more complex environments where we expect this algorithm to shine in the future.

\subsection{Environment and Experiment Setup}
%\vspace{-0.1cm}
The agent’s hull is supported by two legs (the agent appears on the left edge of figure \ref{fig:typical_landscape}).
%\jc{I prefer the standard format of `Fig. 1, left', but we need to be consistent}
The hips and knees of each leg are controlled by two motor joints, creating an action space of four dimensions. The agent has ten LIDAR rangefinders for perceiving obstacles and terrain, whose measurements are included in the state space. Another 14 state variables include hull angle, hull angular velocity, horizontal and vertical speeds, positions of joints and joint angular velocities, and whether legs touch the ground \cite{brockman}. 
%\ks{What do you mean there are no coordinates?  This sentence has no context.  No coordinates for what?}\rw{deleted ``There are no coordinates.''}

Guided by its sense of the outside world through LIDAR and its internal sensors, the agent is required to navigate, within a time limit and without falling over, across an environment of a generated terrain that consists of one or more types of obstacles. These can include stumps, gaps,
%\jc{ I don't think most people know what this word means. I think we should just call these holes throughout. also, these are not actually gaps according to the def: https://goo.gl/3tb15y b/c not hidden}, 
and stairs on a surface with a variable degree of roughness, as illustrated in Figure \ref{fig:typical_landscape}. 
\begin{figure}
  \centering
  \includegraphics[width=.99\linewidth]{figures/new_banner2}
\caption{\textbf{A landscape from the Bipedal Walker environment}. Possible obstacles include stumps, gaps, stairs, and surfaces with different amounts of roughness.}
  \label{fig:typical_landscape}
\end{figure}
As described in the following equation, reward is given for moving forward, with almost no constraints on agents' behaviors other than that they are encouraged to keep their hulls straight and minimize motor torque. If the agent falls, the reward is $-100$. 
%\ks{Language is imprecise: What is 300+ from?  Exactly how is reward computed?}\ks{How is that encouraged?}\rw{addressed}
%\begin{gather*}\label{eq:reward}
\[
    \text{Reward per step}= 
\begin{cases}
    -100,& \text{if robot falls }\\
    130 \times \Delta x-5\times\Delta \text{hull\_angle} - 0.00035\times \text{applied\_torque},& \text{otherwise.}
\end{cases}
\]
%\end{gather*}
The episode immediately terminates when the time limit (2,000 time steps) is reached, when the agent falls, or when it completes the course. In this work, we define an environment as \emph{solved} when it both reaches the far end of the environment and obtains a score of 230 or above.
Based on our observations, the score of 230 ensures that the walker is reasonably efficient.

%\ks{Is the neural network setup based on a previously available implementation or paper?  If so, that should be stated and the original should be cited.}\rw{added}
Following the architecture of controllers for the same domain by \citet{davidha:arxiv18},
all controllers in the experiments are implemented as neural networks with 3 fully-connected layers with $tanh$ activation  functions. The controller has 24 inputs and 4 outputs, all bounded between \text{-}1 and 1, with 2 hidden layers of 40 units each. Each ES step (Algorithm \ref{alg:es_step}) has a population size of 512 (i.e.\ the number of parameters sampled to create the weighted-average for one  gradient step is 512) and and updates weights through the Adam optimizer \cite{kingma2014adam}.
%\jl{suggest: and updates weights through the Adam optimizer}. 
The learning rate is initially set to $0.01$, and decays to $0.001$ with a decay factor of $0.9999$ per step. The noise standard deviation for ES
%\jl{I think this is different terminology than used when introducing ES in the background -- we should be consistent; also the claim there was that the variance was fixed, not adaptive as it is here.} 
is initially set to $0.1$, and decays to $0.01$ with a decay factor of $0.999$ per step. When any environment-agent pair accepts a transfer or when a child environment-agent pair is first created, we reset the state of its Adam optimizer, and set the learning rate and noise standard deviation to their initial values, respectively. 
%\ks{Aren't there any other ES parameters?}\rw{added here and in section \ref{evolution strategies details}}

\subsection{Environment Encoding and Mutation}

Intuitively, the system should begin with a single flat and featureless environment 
%\JC{ I like providing this intuitive explanation, but shouldn't you go first instead of after the morass of details?} \jc{ will it really start out as flat? We just explained that we sample uniformly the roughness of the surface} \jc{ since flatness is an independent dimension of variation should we also say that it's flat and featureless in the sense that it does not have stumps holes Etc?} 
whose paired policy can be optimized easily (to walk on flat ground).  
From there, new environments will continue to be generated from their predecessors, while their paired policies are simultaneously optimized.
%\jc{ this second Clause is confusing. I can't fix it because I don't know exactly what we're trying to say here}  
The hope is that a wide variety of control strategies and skill sets will allow the completion of an ever-expanding set of increasingly complex obstacle courses, all in a single run.  

To enable such a progression, we adopt a simple encoding 
%\jl{will readers understand the difference between direct and indirect encodings?}\jc{ I agree with Joel. They will not. This is yet another example of us using the language of evolutionary algorithms without introducing and defining it. many of our readers are from the reinforcement learning community at are not  familiar with  these terms, and even if we were writing this for the evolutionary algorithms Community we should still Define these terms for those in that community that do not know these terms} 
to represent the search space of possible environments. 
As enumerated in Table \ref{env-mutation-table}, there are five types of obstacles that can be distributed throughout the environment: (1) stump height, (2) gap width, (3) step height, (4) step number (i.e.\ number of stairs), and (5) surface roughness. 
%\jl{Is step height etc. really a type of obstacle, or a parameterization of a type of a basic obstacle, like a stump or a step or a gap?}\jc{the same question could be asked for roughness of surface}\rw{how about we call them ``five types of obstacle parameters''} 
%\ks{You call these "genomes" but they are "genes."}\rw{removed the incorrect or unnecessary usage of genome}
Three of these obstacles, e.g.\ stump height, are encoded as a pair of parameters (or \emph{genes}) that form an interval from which the actual value for each instance of that type of obstacle in a given environment is uniformly sampled.  As will be described below, in some experiments, some obstacle types are intentionally omitted, allowing us to restrict the experiment to certain types of obstacles. 
%\ks{what is an inactive value? zero?}\rw{yes. rewritten}    
%The first value an obstacle parameter takes when it goes from being absent 
%\ks{not clear if "being absent" refers to the case just explained when the whole experiment lacks an obstacle, or whether  you are implying there is some special state where an obstacle type is "absent" that can be switched back to "present"}\rw{modified.}
%to being present in an environment is called its \emph{initial value} (initial values are shown in Table \ref{env-mutation-table}). 
%\ks{clarify how it works}\rw{rewritten} 
When selected to mutate for the first time, obstacle parameters are initialized to the corresponding initial values shown in Table \ref{env-mutation-table}. 
%\jl{This previous sentence is a bit confusing/imprecise -- if you mutate the step height parameter for the first time, it also initializes the step number, right? And it seems to muddle the issue just explained that there is no 'step height' parameter (singular), but two paired parameters.} 
For subsequent mutations, an obstacle parameter takes a \emph{mutation step} (whose magnitude is given in Table \ref{env-mutation-table}), and either adds or subtracts the step value from its current value.
%\jc{ this language implies there is a singular value, but some types of obstacles there are two values}. 
Note that for surface roughness, both the initial value and every subsequent mutation step are sampled uniformly from $(0, 0.6)$ because the impact of roughness on the agent was found to be significant sometimes even from very small differences. 
%\jc{ why?} 
The value of any given parameter cannot exceed its \emph{maximum value}. %A typical terrain is generated by alternating between rough surface with random length whose roughness is determined by the Roughness genome and one of the three types of obstacles whose values are determined by the value of the corresponding genome(s) if active. 

Because the parameters of an environment define a distribution, the actual environment sampled from that distribution is the result of a random \emph{seed}.  This seed value is stored with each environment so that environments can be be reproduced precisely, ensuring repeatability.  (The population of many environments and the mutation of environments over time still means that training overall does not occur on only one deterministic environment.)
With this encoding, all possible environments can be uniquely defined by the values for each obstacle type in addition to the seed that is kept with the environment. 

Any environments that satisfy the reproduction eligibility condition in Line 7 of Algorithm \ref{alg:mutate_niches} (which is given in Supplemental Information) are allowed to mutate to generate a child environment. In our experiments, this condition is that the paired agent of the environment achieves a reward of 200 or above,  which generally indicates that the agent is capable of reaching the end of the terrain (though slightly below the full success criterion of 230). 
To construct \texttt{child\_list} in Line 15 in Algorithm \ref{alg:mutate_niches}, the set of eligible parents is sampled uniformly to choose a parent, which is then mutated to form a child that is added to the list.  This process is repeated until \texttt{child\_list} reaches \texttt{max\_children}, which is 512 in our experiments.  Each child is generated from its parent by independently picking and mutating some or all the available parameters of the parent environment and then choosing a new random seed.
%To construct \texttt{child\_list} in Line 15 in Algorithm \ref{alg:mutate_niches}, each parent environment first generates a number of (here, \todo{X}) child environments 
%by independently picking and mutating some or all the available parameters of the parent environment and then choosing an \emph{environment seed}, which ensures that any rollout of the same environment yields exactly the same terrain\jc{ unclear. the reader will not know what a roll-out is and how the terrain of an environment can differ from the environment itself. I suggest we describe  set each environment represents a distribution of  possible instances of that environment type and that the random seed produces the same instance from that distribution each time}. Each of the newly created child environments are then evaluated against the parent environment's current paired agent.\jc{ the language hear sounds backwards I would think that you evaluate agents in an environment not evaluate environments against agents} 
%\ks{but isn't there a minimal criterion or something like that?  what if it's too hard?}\rw{corrected}. 
The \emph{minimal criterion} (MC) $50 \leq E^{\textrm{child}}(\theta^{\textrm{child}}) \leq 300$ 
%\ks{shouldn't F in this mathematical expression be E?}\rw{corrected. here and below}
then filters out child environments that appear too challenging or too trivial for the current level of capability of agents. In case the number of child environments that satisfy the MC is more than the maximum number of children admitted per reproduction, those with lower $E^{\textrm{child}}(\theta^{\textrm{child}})$ are admitted until the cap is reached. 
%\jc{ Vague. How?  do we take them in that order until we reach the cap?}. 

\begin{table}[H]
\begin{center}
\begin{small}
\begin{sc}
\begin{tabular}{lccccc}
\hline
\begin{tabular}{@{}c@{}}Obstacle \\ Type\end{tabular} & \begin{tabular}{@{}c@{}}Stump \\ Height\end{tabular} & \begin{tabular}{@{}c@{}}Gap \\ Width\end{tabular} & \begin{tabular}{@{}c@{}}Step \\ Height\end{tabular} & \begin{tabular}{@{}c@{}}Step \\ Number\end{tabular} & Roughness\\


\hline
Initial Value & $(0.0, 0.4)$ & $(0.0, 0.8)$ &  $(0.0, 0.4)$ & $1$ & uniform$(0, 0.6)$\\
Mutation Step & $0.2$ & $0.4$ & $0.2$ & $1$ & uniform$(0, 0.6)$\\
Max Value & $(5.0, 5.0)$ & $(10.0, 10.0)$ & $(5.0, 5.0)$ & $9$ & $10.0$\\
\hline
\end{tabular}
\end{sc}
\end{small}
\end{center}
\caption{\textbf{Environmental parameters (genes) for the Bipedal Walker problem.} 
%\ks{Rui: Somehow the table now goes past the right margin.  Not sure what happened.}\rw{fixed. commented this out.} 
}\label{env-mutation-table}
\vskip -0.1in
\end{table}

%\ks{added}


\subsection{Results} 

%\ks{Also missing in this section: We should give a sense of how much computation (e.g. how many iterations of ES) it takes to get to certain points.  For example, how long does it take to get to the targets shown in figure 4?}\rw{worker information added when parallel implementation is described. More included in texts for Figure \ref{fig:control_type_1} in this section.}\ks{I'm not seeing where text for figure 4 explains the *amount* of computation -- did I miss it?  What I am suggesting is that the reader would like to know roughly how many steps it takes to run a whole run, and how many to get to challenging red pentagons.  We don't need a big table of every run length; is there some kind of brief summary you can give instead to characterize the general computation requirements to get to where we did?}\rw{in the paragraph above Table \ref{challenge-table} i reported that each single run of POET is up to 25,200 iterations. if i plan to add some information about how many to get to challenging red pentagons: for example, averagely how many POET iterations to take to solve level 1 envs, level 2 envs, level 3 envs, my questions is which number to report. for example, a difficult level 3 environment is first created at Iteration 8000,  "solved" at Iteration 9786, retired at Iteration 11000. Should I report 9786? or 1786? or 3000? I suggest we report that each environment once admitted, will all stay in system for 3000 POET iterations. Then report average POET iterations actually spent on : in this case 1786. Please see the proposed sentences at the paragraph before Table  \ref{challenge-table}}\ks{agreed, let's go with your proposed approach of reporting the iterations spent (e.g. 1,786); also remember to use commas in numbers of 3 digits long}\rw{agree. will update the number in the proposed place labelled as TODO}

An important motivating hypothesis for POET is that the stepping stones that lead to solutions to very challenging environments are more likely to be found through a divergent, open-ended process than through a direct attempt to
optimize in the challenging environment.  
Indeed, if we take a particular environment evolved in POET and attempt to run ES (the same optimization algorithm used in POET) on that specific environment from scratch, the result is often premature convergence to degenerate behavior.
%\jc{should we also mention that we'll do/beat the obvious curricula control where we increase complexity towards the target?}\ks{This point is too nuanced to fit in here -- we do not simply beat the control in all cases.  We only beat it when the target environment is sufficiently complex.  Explaining that here would be a distraction from the simple initial point, and we'll get to the curricula later anyway.} 
The first set of experiments focus on this phenomenon by running POET with only \emph{one} obstacle type enabled (e.g.\ gaps in one case, roughness in another, and stumps in another). %\ks{important to be transparent on what kind of experiment was run to get reported results}\rw{agree}.  
That way, we can see that even with a single obstacle type, the challenges generated by POET (which POET solves) are too much for ES on its own. 

Figure \ref{fig:comparison_naive_control} shows example results that illustrate this principle.  Three challenging environments that POET created are shown, with wide gaps, rugged hillsides, and high stumps, respectively. To demonstrate that these environments are challenging, we ran ES to directly optimize randomly initialized agents for them. More specifically, for each of the three environments, we ran ES five times with different initialization and random seeds up to 16,000 ES steps (which is twice as long as \citet{davidha:arxiv18} gives agents in the same domain with ES). Such ES-optimized agents are consistently stuck 
%\ks{How consistently? Can we quantify it at all?  How many times did why try to solve each with ES?  How many times did it fail?  Does it always fail or just usually?  How usually?  Is there a numerical statistic we can give to characterize this failure more objectively?  The reader here will want to see evidence beyond just the words ``consistently stuck.''}\rw{please see the proposed text.}  
at local minima in these environment. 
The maximum scores out of the five ES runs for the environments illustrated in Figure \ref{fig:large_gap}, \ref{fig:hillside}, and \ref{fig:high_stump} are 17.9, 39.6, and 13.6, respectively, far below the success threshold of 230 that POET exceeds in each case.  
In effect, to obtain positive scores, these agents learn to move forward, but also to freeze before challenging obstacles, which helps them avoid the penalty of $-100$ for falling.  This behavior is a local optimum: the agents could in principle learn to 
overcome the obstacles, but instead converge on playing it safe by not moving.
%\jl{Should there be explicit language here describing this as a local optimum? Here and then in the next paragraph it sounds almost like these are generally risk-averse agents, intentionally not taking risks, whereas the truth is more like they've converged to a local optimum that narrowly recognizes some particular obstacle and stops.}\jc{ I agree with joel} 
Note that ES had previously been shown to be a competent approach \cite{dha_esblog} on the original Bipedal Walker Hardcore environment in OpenAI Gym, which, as illustrated later in Table \ref{challenge-table}, consists of similar, but much less difficult obstacles.
The implication is that the challenges generated and solved by POET are significantly harder and ES alone is unable to solve them.
%\jc{ this leaves are hypothesis for why the results are different a little big and unstated and implied. Let's be more explicit. do we think the reason is because these are harder challenges?}

%\ks{Should we say something about David Ha's ES experiments in the same domain here?  That is, the point would be to emphasize that ES is a competent (maybe the best) approach for this domain, which gives emphasis to the impressiveness of our improvements beyond what raw ES can do on its own.}\rw{agree. done. David Ha tried ES on OpenAI Gym's original bipedal walker hardcore, which is much less difficult, as illustrated in Table \ref{challenge-table}. I am citing that blog. }

\begin{figure}
    \centering
    \begin{subfigure}[b]{1.0\textwidth}
        \includegraphics[width=\textwidth]{figures/oe-gap2}
        \caption{ES-/POET-generated agents attempting gaps}
        \label{fig:large_gap}
    \end{subfigure}
    \begin{subfigure}[b]{0.49\textwidth}
        \includegraphics[width=\textwidth]{figures/roughness_comp3}
        \caption{ES-/POET-generated agents on rough surfaces}
        \label{fig:hillside}
    \end{subfigure}
    ~ %add desired spacing between images, e. g. ~, \quad, \qquad, \hfill etc. 
      %(or a blank line to force the subfigure onto a new line)
    \begin{subfigure}[b]{0.49\textwidth}
        \includegraphics[width=\textwidth]{figures/stump_comp3}
        \caption{ES-/POET-generated agents attempting stumps}
        \label{fig:high_stump}
    \end{subfigure}
    \caption{\textbf{POET creates challenging environments and corresponding solutions that cannot be obtained by optimizing randomly initialized agents by ES.} As illustrated in the top row of (a) and the left panels of (b) and (c), agents directly optimized by ES converge on degenerate behaviors that give up early in the course. In contrast, POET not only creates these challenging environments, but also learns agents that overcome the obstacles to navigate effectively, as shown in the bottom row of (a) and right panels of (b) and (c). 
    %\ks{The graphic in this figure should be easier to understand without having to read the caption.  Maybe groupings could by made clear by drawing boxes aroiund them.  Also, annotations like "ES from scratch" and "POET-generated" that help us know what is what would make it much clearer much faster.  Note also: subscaptions (a) (b) and (c) should have very brief text next to them describing what they are about. }\rw{done}
    } \label{fig:comparison_naive_control}
\end{figure}

For example, in Figure \ref{fig:large_gap}, the sequence in the top row illustrates the agent choosing to stop to avoid jumping over a dangerously wide gap: it slowly sticks one of its feet out to reach the bottom of the first wide gap and then maintains its balance without any further movement until the time limit of the episode is reached.
Videos of this and other agents from this paper can be found at \url{eng.uber.com/poet-open-ended-deep-learning}.
%\jc{ this would be a good place to say that videos of these and other Agents from this paper can be viewed at and include a URL} 
In contrast, POET not only \emph{creates} such challenging environments, but also in this case learns a clever behavior to overcome wide gaps and reach the finish line (second row of Figure \ref{fig:large_gap}). 
%\jc{ we provided the numbers for ES, but also need to provide the results numbers for poet and include a  box plot for each of the environments and do mann-whitney statistical test to show that our results are statistically better}\ks{it's not clear what kind of comparison you have in mind here because ES and POET are not directly comparable statistically here.  The reason is that POET does not try to reach a pre-conceived target behavior the way ES does, so you cannot run POET n times and say, ``this is the score it achieves on average for this target.''  This subtlety is one of the reasons that the proper assessment of open-ended systems is a genuinely complex and controversial topic.  open-ended systems by definition cannot be asked to reproduce a defined target result, so you can't see how consistent they are in that way.  much of the explication of results here is an attempt to navigate around this roadblock.}

Figure \ref{fig:challenging_envs} illustrates two more interesting examples of challenging environments and the highly adapted 
%\jl{I think hyphenation is grammatically correct here?}\jc{ I have looked this up before and do not think it is necessary because of the `ly'} 
capabilities of their paired agents.   In the first, the agent navigates a very steep staircase of oversized steps, which ES alone cannot solve.  In the second, an experiment in heterogeneous environments (with both gaps and stumps available) yields a single obstacle course of varying gap widths and stump heights where
the POET agent succeeds and ES fails. 
%\jc{ again we need to provide the actual data and box plots for both es and poet and statistical tests. we can't just report anecdotal stories}  
%\ks{Rui: can we at minimum report the max performance of ES on the environments shown in figure 3 as well?}\rw{done, todo cleared} 
The maximum scores out of the five ES runs for the environments illustrated in Figure \ref{fig:downstairs_hero} and \ref{fig:mixed_up} are 24.0 and 19.2, respectively, again far below the success theshold for POET of 230. Figure \ref{fig:compare-poet-es} illustrates the dramatic failure of ES to match the performance of POET in these relatively simple environments that POET generated and solved.
%Videos of examples of environments POET creates and their paired agents navigating in these environments can be found at \url{eng.uber.com/poet-open-ended-deep-learning}.\jc{ move this up to the more natural place I pointed out in the last paragraph}

\begin{figure}
    \centering
    \begin{subfigure}[b]{0.49\textwidth}
        \includegraphics[width=\textwidth]{figures/ds_hero2}
        \caption{A downstairs course with a very large step height}
        \label{fig:downstairs_hero}
    \end{subfigure}
    ~ %add desired spacing between images, e. g. ~, \quad, \qquad, \hfill etc. 
      %(or a blank line to force the subfigure onto a new line)
    \begin{subfigure}[b]{0.49\textwidth}
        \includegraphics[width=\textwidth]{figures/mixed_env2}
        \caption{A rough walking course with gaps and stumps}
        \label{fig:mixed_up}
    \end{subfigure}
    \caption{\textbf{Agents exhibit specialized skills in challenging environments created by POET}.  These include the agent navigating (a) very large steps and (b) a rough walking course with mixed gaps of small and large widths, and stumps of various heights. 
    %\ks{the right-hand side of panel (a) looks confusing and weird.  maybe just show it at one scale (in this figure) because the reader here is only looking at still snapshots and has no idea why (or if) the agents is being displayed at two separate scales.}\rw{Done. fixed (a). replaced (b). add subcaption}
    } \label{fig:challenging_envs}
\end{figure}


\begin{figure}
  \centering
  \includegraphics[width=.75\linewidth]{figures/es_vs_poet2}
\caption{\textbf{ES on its own cannot reproduce results achieved by POET.} Each column compares a box plot representing five ES runs to the performance achieved by POET in the single-obstacle-type problems (except for the last column) generated by POET.  Recall that a score of 230 (the dotted line) is the threshold for success.  
Though they are hard to see in these plots because of the narrow ES distributions, the red lines are medians, which are surrounded by a box from the first quartile to the third, which are further surrounded by lines touching the minimum and maximum. 
As the plot illustrates, ES on its own definitively fails in every case to come close to the challenges set and solved by POET.}
  \label{fig:compare-poet-es}
\end{figure}

%\todo{work in this argument}
An interesting aspect of open-ended algorithms is that
statistical comparisons can be challenging 
because such algorithms are not trying to perform well on a specific preconceived target, and thus some amount of post-hoc analysis is often necessary.  That property is what
motivates an analysis driven by choosing environments produced and solved by POET, and then checking whether direct optimization can solve these environments.
However, from a statistical comparison standpoint it raises the challenge that because each environment created by POET is unique, we only have one score for POET to compare against a set of direct-optimization ES runs. For that reason, we use a \emph{single-sample t-test} that assumes the POET distribution is centered on the single POET score and measures whether the direct optimization distribution’s mean is statistically significantly lower. 
As the very poor results for ES would suggest, based on single-sample t-tests for all five environments in figures \ref{fig:comparison_naive_control} and \ref{fig:challenging_envs} that are created by POET, the scores of agents optimized by ES alone are very unlikely to be from a distribution near the POET-level solutions
($p < 0.01$). 
%\ks{We should scrutinize claims we make based on single-sample t-tests because this is an unconventional way to analyze results.  Does this kind of statement really add value, or does it actually raise unnecessary questions (i.e. about the statistical test and its usage)?}

One interpretation of POET's ability to create agents that can solve challenging problems is that it is in effect an automatic curriculum builder. Building a proper curriculum is critical for learning to master tasks that are challenging to learn from scratch due to a lack of informative gradient. However, building an effective curriculum given a target task is itself often a major challenge. Because newer environments in POET are created through mutations of older environments and because POET only accepts new environments that are not too easy not too hard for current agents, POET implicitly builds a curriculum for learning each environment it creates.  The overall effect is that it is building many overlapping curricula simultaneously, and continually checking whether skills learned in one branch might transfer to another.

A natural question then is whether the environments created and solved by POET can also be solved by an explicit, \emph{direct-path curriculum-building control} algorithm.  To test this approach, we first collect a sample of environments generated and solved by POET, and then apply the direct-path control to each one separately to see if it can reach the same capabilities on its own. 
%\jc{ the use of the word "the "  is not right here because we have two controls vanilla ES and this curriculum Builder.  also every paper has controlled so this name does not really describe the algorithm period we should come up with a better name.  how about straight line curriculum Builder?  because it takes the straight line between the starting point and the target objective.}, 
%works by first setting a target environment, e.g.\ a challenging one created and solved by POET\jc{ this makes it sound like the control algorithm sets the Target environment. we should just say what we did which is that we manually a few challenging environments produced by poet and then use them as the target environment control}. 
In this control, the agent is progressively trained on a sequence of environments of increasing difficulty that move towards the target environment. 
%\jl{Is there any cite for a similar approach to curriculum building in the literature, to give it more credibility? The parkour paper had some kind of curriculum, or there's sebastian's work with PCG that I think trained from easy to hard, etc.} 
%\ks{Rui: Joel is right we should cite something in the next sentence below to give us credibility for our control through precedent.  One clear example is Faustino Gomez's 1997 paper, Incremental Evolution of Complex General Behavior (for which I added the cite), which is at http://www.cs.utexas.edu/users/nn/downloads/papers/gomez.adaptive-behavior.pdf.  But also Rui, can you check for the other ones Joel mentioned too, e.g. the parkour one and maybe Sebastian's?}\rw{Done. removed todo. I checked and they are relevant and should be cited, so I added parkour, Sebastian's PCG paper and two extra important cites about curriculum mentioned in the first two paper.}
This kind of incremental curriculum is intuitive and variants of it appear in the literature when a task is too hard to learn directly \cite{gomez:ab97, bengio2009curriculum, karpathy2012curriculum,heess2017emergence, justesen2017procedural}.
The sequence of environments start with an environment of only flat ground without any obstacles (which is easy enough for any randomly initialized agent to quickly learn to complete). Then each of the subsequent environments are constructed by slightly modifying the current environment. More specifically, to get a new environment, each obstacle parameter of the current environment has an equal chance of staying the same value or \emph{increasing} by the corresponding mutation step value in Table \ref{env-mutation-table} until that obstacle parameter reaches that of the target environment. 
%\ks{confusing: *mutation* is random -- is this process random?  I thought it moves in a non-random direction towards the target?} \rw{rewritten} 
%\ks{Still a little unclear: "increasing at least one" doesn't tell me how we decide which one(s) to increase, or why it would choose one to increase and not another?}\rw{rewritten.}

In this direct-path curriculum-building control, the agent moves from its current environment to the next one when the agent's score in the current environment reaches the reproduction eligibility threshold for POET, i.e.\ the same condition for when an environment reproduces in POET. The control algorithm optimizes the agent with ES (just as in POET). It stops when the target environment is reached and solved, or when a computational budget is exhausted. To be fair, each run of the control algorithm is given the same computational budget (measured in total number of ES steps) spent by POET to solve the environment,
which includes all the ES steps taken in the entire sequence of environments along the direct line of ancestors (taking into account transfers) leading to the target. 
This experiment is interesting because if the direct-path control cannot reach the same level as the targets, it means that a direct-path curriculum alone is not sufficient to produce the behaviors discovered by POET, supporting the importance of multiple simultaneous paths and transfers between them.
%\rw{propose to add: which includes independent ES optimization steps for the environment and all its predecessor environments as well as all calculations in transfer experiments with the environment and all its predecessor environments as target.}\jc{ this explanation is confusing. we need to add these details but we need to make the writing clearer} \ks{might want to clarify what we mean by budget spent by POET -- do we mean just in that environment or do we mean stretching back to earlier predecessor environments, or do we mean the budget of the entire POET run including solving other environments?}\rw{add some texts above to clarify this.} \jc{I think instead we should discuss in email whether to adopt the language of an archive only (with active and frozen environments) or a population and an archive, or something else. then we should implement that solution.}

The direct-path curriculum-building control is tested against a set of environments generated and solved by POET that encompass a range of difficulties. 
For these experiments, the environments have three obstacle types enabled: gaps, roughness, and stumps.
Unlike when comparing to ES alone, he we allow POET to generate environments with multiple obstacles at the same time because the curriculum-based approach should in principle be more powerful than ES alone.
To provide a principled framework for choosing the set of generated environments to analyze, they are classified into three difficulty levels. 
The difficulty level of an environment is based on how many conditions it satisfies out of the three listed in Table \ref{challenge-table}.
In particular, a \emph{challenging} environment satisfies one of the three conditions; a \emph{very challenging} environment satisfies two of the three; and satisfying all three makes an environment \emph{extremely challenging}.
It is important to note that these conditions all merit the word ``challenging'' because they all are much more demanding than the corresponding values used in the original Hardcore version of Bipedal Walker in OpenAI Gym \cite{OpenAI_bipedal} (denoted as \emph{reference} in Table \ref{challenge-table}, which was also used by \citet{davidha:arxiv18}.  More specifically, they are 1.2, 2.0, and 4.5 times the corresponding reference values from the original OpenAI Gym environment, respectively. 

\begin{table}[H]
\begin{center}
\begin{small}
\begin{sc}
\begin{tabular}{lccccc}
\hline
 & Top Value & Top Value &  \\
 & in range of &  in range of & Roughness \\
 & Stump Height & Gap Width & \\
\hline
This Work & $\geq2.4$ & $\geq6.0$ &  $\geq4.5$\\
Reference & $2.0$ & $3.0$ & $1.0$\\
\hline
\end{tabular}
\end{sc}
\end{small}
\end{center}
\caption{\textbf{Difficulty level criteria.} The difficulty level of an environment is based on how many conditions it satisfies out of the three listed here. The \emph{reference} in the second row shows the corresponding top-of-range values used in the original Hardcore version of Bipedal Walker in OpenAI Gym \cite{OpenAI_bipedal}. The  difficulty of the environments produced and solved by POET are much higher than the reference values.}\label{challenge-table}
\vskip -0.1in
\end{table}

In this experiment, each of three runs of POET takes up to 25,200 POET iterations with a population size of 20 active environments, while the number of sample points for each ES step is 512.
These runs each take about 10 days to complete on 256 CPU cores.
%Each environment, once admitted, stays active in the population for 3,000 POET iterations.
%\jc{ I believe this is the first time we have mentioned this detail. is this true for all of poet and all of the experiments in this paper or just for these experiments? shouldn't this be in methods?} 
%\ks{Rui: Jeff asks a good question, and also, it's not clear how we guarantee that each environment stays in the population for a fixed number of iterations?  Doesn't that depend on how many child environments are admitted over time?  I think maybe we should even leave out this line because it's confusing and could lean down a complicated tangent discussion, but not sure the answer myself.}\rw{Thank you and Jeff for catching this. I agree we should leave out this line. 3000 is very rough oversimplified estimation that i simply multiply the number of niches with interval of checking if we can add the children. That was too rough. As you pointed out, it really depend on how parents are quickly becoming mature enough to reproduce and how child environments are admitted over time. }
The mean (with 95\% confidence intervals) POET iterations spent on solving challenging, very challenging, and extremely challenging environments starting from the iteration when they were first created are $638\pm133$, $1$,$180\pm343$, and $2$,$178\pm368$, respectively.
Here, one POET iteration refers to 
creating new environments, optimizing current paired agents, and attempting transfers (i.e.\ lines 4-22 in Algorithm \ref{alg:poet}). 
%\jc{nice to have these numbers. but are they within one run (over many environements), across runs, or both? also, do we say how many runs of poet we did?}
The more challenging environments clearly take more effort to solve.
%\jc{add confidence intervals  or some other metric of spread}\rw{Done. cleared todo}
%\ks{Rui: Is the mean POET iterations just mentioned computed starting from when the environment was first created?  Or does it encompass even iteratons spent on predecessor environments? I'd like to clarify this point in the text above.} \rw{it is computed starting from when the environment was first created, not include iterations spent on predecessors.}
%\ks{The above addition seems good, but one confusing details is that the "population" size (which is the number of environments at one time) could be confused for the number of pseudo-offspring for ES steps.  Perhaps that should also be stated above so everything is completely clear.}\rw{agree. added above}



\begin{figure}
  \centering
  \includegraphics[width=.83\linewidth]{figures/figure4_0126_3}
\caption{\textbf{POET versus direct-path curriculum-building controls.} Each rose plot depicts one environment that POET created and solved (red pentagon). For each, the five blue pentagons indicate what happens in control runs when the red pentagon is the target. 
Each blue pentagon is the closest-to-target environment solved by one of the five independent runs of the control algorithm. The five vertices of each pentagon indicate roughness (\emph{Roughness}), the bottom and top values of the range of the gap width of all the gaps (\emph{Gap\_lower} and \emph{Gap\_upper}), and the bottom and top values for the height of stumps  (\emph{Stump\_lower} and \emph{Stump\_upper}) in the given solved environment. The value after MAX in the key is the maximum value at the outermost circle for each type of obstacle. Each column contains sample solved environments from a single independent run of POET.
%\jc{ why do we repeat the legends on each axis of the Rose three times when it's the same information?}\jc{ why we calling it GAP A&B instead of Max and min?}\ks{Rui:good questions from Jeff; can we have the legend only on the plot in the upper left, and use max and min instead of a and b?}\ks{Rui: also, can we label the *rows* as well, such that we label the difficulty level 1, level, and level 3 of the rows?}\rw{New Figure DONE.}
%\ks{Rui: ``Extreme Challenging'' should be ``Extremely Challenging'' (missing the -ly of the adverb extremely).}\rw{fixed. Done. Comment out this.}
}
  \label{fig:control_type_1}
\end{figure}



Figure \ref{fig:control_type_1} compares the POET environments and the direct-path curriculum-building control algorithm through a series of \emph{rose plots}, wherein each rose plot compares the configuration of an environment solved by POET (red pentagons) with the closest that the direct-path control could come to that configuration (blue pentagons).  For each red pentagon there are five such blue pentagons, each representing one of five separate attempts by the control to achieve the red pentagon target.
The five vertices of each pentagon indicate roughness, the lower and upper bounds of the range of the gap width, and those of the stump height, respectively.
Each column in figure \ref{fig:control_type_1} consists of six representative samples (the red pentagons) of environments that a single run of POET up to 25,200 iterations 
%\ks{may want to remind reader/define what iteration means}\rw{added as the next sentence.} 
created and solved.  
As the rows descend from top to bottom, the difficulty level decreases.  In particular, the top two rows are extremely challenging targets, the next two are very challenging, and the bottom two are challenging (all are randomly sampled from targets generated and solved at each difficulty level in each run).  
Qualitatively, it is clear from figure \ref{fig:control_type_1} that attempts by the direct-path control consistently fail to reach the same level as POET-generated (and solved) levels.  
%In each column of Figure \ref{fig:control_type_1}, the top two 
%red pentagons\jc{these difficulty levels should be axes labels on the rows in the figure.  for example we could call them challenging, very challenging and extreme challenging}\ks{Rui: I agree with Jeff, made a similar comment elsewhere.  We can use Jeff's suggested wording for these labels too}\rw{DONE} depict two environments sampled \jc{ sampled how? randomly ?by hand?}\ks{Rui: I believe these are randomly sampled (and that would be best); can you confirm that?}\rw{Yes, Confirmed.}  from those that satisfy all three conditions, the middle two red pentagons depict two environments sampled from those satisfying only two conditions, and bottom two red pentagons depict two environments sampled from those satisfying only one condition.
%The five blue pentagons (each representing a separate, randomly initialized run of the single-path curriculum control) in each rose plot then show what happens with the control if the red pentagon is set as the target.  In particular, the blue pentagons are the performance of the resultant agent in the closest-to-target environments that five independent runs of the control algorithm can solve. 

%HERE-KEN

To more precisely quantify these results, we define the normalized distance between any two environments, $E_{A}$ and $E_{B}$ as
$ \frac{1}{\beta}\vert\vert{\frac{e(E_{A}) - e( E_{B})}{e(E_{\textrm{Max}})}}\vert\vert_{2}$, where $e(E)$ is the genetic encoding vector of $E$, and $E_{\textrm{Max}}$ is a hypothetical environment (for normalization purpose) with all genetic encoding values are maxed out. (The maximum values are $\textrm{roughness} = 8, \textrm{Gap\_lower}=\textrm{Gap\_upper} = 8, \textrm{Stump\_lower} = \textrm{Stump\_upper} = 3$.)
%\ks{Rui: as Jeff suggested elsewhere, can we replace a and b with min/max or max/min, whichever is correct?  note that when subscripts have multiple letters they will need brackets around them in the latex}\rw{we decided to use lower and upper} 
The constant $\beta = \sqrt{5}$ is simply for normalization so that the distance is normalized to $1$ when $E_{A} = E_{\textrm{Max}}$ and $E_{B}$ is a purely flat environment (roughness zero) without any obstacles, which therefore contains an all-zero genetic encoding vector. 
%\jc{ did we say somewhere before that the flat environment is represented by a roughness of 0 as this sentence implies?}\ks{Something doesn't make sense here: if the distance measure is supposed to be between \emph{any} two environments, then why here are we now saying that $E_{B}$ is purely flat?  Can't $E_{B}$ be anything in the expression for distance?  Why or when should it be only flat?} \rw{Modified the text. Yes, $E_{B}$ be anything in the expression for distance. $E_{B}$ is purely flat, together with $E_{A} = E_{\textrm{Max}}$ is for explanation of how the normalization constant $\sqrt{5}$ is calculated.} 
Based on this distance measure, we can calculate the median values and confidence intervals of distances between target environments (created and solved by POET) and the corresponding closest-to-target environments that the control algorithms can solve. 
These values are calculated for all the environments of different challenge levels across all three runs shown in figure \ref{fig:control_type_1}. 

The results, summarized in figure \ref{fig:distance-boxplot}, demonstrate that POET creates and solves environments that the control algorithm fails to solve at very and extremely challenging difficulty levels, while the curriculum-based control algorithm can often (though not always) solve environments at the lowest challenge level. 
One implication of these results is that the direct-path curriculum-building control is valid in the sense that it does perform reasonably well at solving minimally challenging scenarios.  However, the very and extremely challenging environments that POET invents reach significantly beyond what the direct-path curriculum can match.
%validating the efficacy of the control but also showing its limitations relative to POET-style search\jc{ this last Clause is too much to pack into an already long complicated sentence. Break it out to its own sentence and explain it better}.
There are a couple ways to characterize this significance level more formally.   These analyses are based on the \emph{distance between the closest environment reached by the control and the POET-generated target}.  First, we can examine whether that distance is significantly greater the higher the challenge level.  Indeed, Mann-Whitney U tests 
show that the difference between such distances between challenging and very challenging environments is indeed very significant ($p < 0.01$), as is the difference between such distances between very challenging and extremely challenging.  This test gives a quantitative confirmation of the intuition (also supported by figure \ref{fig:distance-boxplot}) that these challenge levels are indeed meaningful, and the higher they go, the more out of reach they become for the control.  Second, 
a single-sample t-test (which was previously used when comparing POET to the ES-alone control)
can also here provide confidence that the distances from the targets are indeed far in an absolute sense: indeed, even at the lowest challenge level (and for all challenge levels), the one-sample t-test measures high significance ($p < 0.01$) between the distribution of distances from the target and the target level.

%\rw{propose to add: 
%\jc{the following paragraph seems to just restate the previous paragraph (but the writing in the previous paragraph is better. I assume we should remove this paragraph right?}
%\todo{remove: Based on Single Sample T-Test, for all three difficulty levels, the distances between targets created and solved by POET and the distribution of what controls could reach, are all very significantly ($p < 0.01$) nontrivial (nonzero). Based on Mann–Whitney U tests, as the difficult level goes up from ``Challenging'' to ``Very Challenging'' and then to ``Extreme Challenging'', there are all statistically very significant increases  ($p < 0.01$) in the distances between targets created and solved by POET and what controls could reach.} 
%\ks{right, thanks, had forgotten to remove when I edited the previous paragraph}

In aggregate, the results from the direct-path curriculum-building control
help to show quantitatively the advantage of POET over conventional curriculum-building. In effect, the ability to follow multiple chains of environments in the same run and transfer skills among them pushes the frontier of skills farther than a single-chain curriculum can push, hinting at the limitations of preconceived curricula in general. 
Thus, while POET cannot guarantee reaching a particular preconceived target,
these results suggest that in some environment spaces POET's unique ability to generate multiple different challenges and solve them may actually still provide a more promising path towards solving some preconceived target challenges (i.e. offering a powerful alternative to more conventional optimization), putting aside the additional interesting benefits of an algorithm that invents, explores, and solves new types of challenges automatically.
%\ks{Jeff: the claim you articulate here doesn't seem to follow.  How can POET be a ``better path towards solving a preconceived target challenge'' if it cannot be directed towards a target?  Nothing in the results implies we could set a specific target for POET.}

\begin{figure}
  \centering
  \includegraphics[width=.85\linewidth]{figures/figure_box_control3}
\caption{\textbf{Distances between targets and results of direct-path curriculum-building control runs.} Here, levels \emph{challenging}, \emph{very challenging}, and \emph{extremely challenging} correspond to the bottom two, middle two and top two rows of red pentagons in Figure \ref{fig:control_type_1}, respectively.
The red lines are medians, which are surrounded by a box from the first quartile to the third, which are further surrounded by lines touching the minimum and maximum.
These results highlight how difficult it becomes for the direct-path control to match challenges invented and solved by POET as the challenge level increases.
}
  \label{fig:distance-boxplot}
\end{figure}

\begin{comment}

\begin{table}[H]
\begin{center}
\begin{small}
\begin{sc}
\begin{tabular}{lccccccc}
\hline

 Difficulty & \begin{tabular}{@{}c@{}}Number of \\ Condition \\ Satisfied\end{tabular}&  \multicolumn{4}{c}{\begin{tabular}{@{}c@{}}Distance \\ Between Target and Controls \end{tabular}} \\
 \cline{3-6}
            &           & Max & Min & Mean & STD  \\
\hline
Level 3 & $3$ & $0.592$ & $0.278$ &$0.375$ &$0.062$\\
Level 2 & $2$ & $0.353$ & $0.028$ & $0.214$ & $0.078$ \\
Level 1 & $1$ & $0.224$ & $0.000$ & $0.081$ & $0.077$\\
\hline
\end{tabular}
\end{sc}
\end{small}
\end{center}
\caption{\textbf{Distances between targets and results of control runs.} Here, Levels ``Extreme Challenging'', ``Very Challenging'', and ``Challenging'' correspond to the top two, middle two and bottom two rows of red pentagons in Figure \ref{fig:control_type_1}, respectively.

}\label{distance-table}
\vskip -0.1in
\end{table}
\end{comment}
%\jc{because we do not know if the data are normally distributed we should use medians and confidence intervals.  it is not appropriate to use mean and standard deviation.}\jc{ we need to do statistics do you say whether these are significantly different than zero. note that are statistics cannot assume normality in this we do test for normality. I think just doing bootstrap confidence intervals of the median and showing that less than 0.01 of them include 0 would work.  we can also show that level two is statistically different than level one and that three is different from 2 and 1.  if we want to get fancy there are tests to show that the three orders of difficulty are in an ordered pattern, ie that 3 is greater than 2 and 2 is greater than 1, but I don't know exactly how to do that so you'd have to research it}\rw{Above is Jeff's comments on old Table of Distance Poet vs control. All addressed: box plots, and significance tests are all added}


%\ks{not clear what "ranges of environments" means here?  the range of stump-height and gap width in a single run?  or the range of something (what?) across multiple runs?}
%\ks{there is a lack of overall quantitative data.  how many runs of POET were there total?  How many or how diverse are its products in general?  Given a method of sampling environments from a single run, how close does the control come to matching their performance on average?  What is the standard deviation?  Is it significant?  The results so far are convincing in an anecdotal sense, but a general statement would help to avoid the appearance of cherry-picking.}
%\rw{modified or added table/texts/plots to address the above comments about Figure 4 and narratives.} 

%\ks{do we have any references for papers that use "conventional" curricula?  We should be referencing them somewhere in this section given we talk about them so much.}\rw{So far I have not found any references.}\ks{still an open issue}\ks{resolved above}

A fundamental problem of a pre-conceived direct-path curriculum (like the control algorithm above) 
is the potential lack of necessary \emph{stepping stones}. In particular, skills learned in one environment can be useful and critical for learning in another environment. Because there is no way to predict where and when stepping stones emerge, the need arises to conduct transfer experiments (which POET implements) from differing environments or problems \cite{nguyen2016innovative}. In conventional curriculum-building algorithms, the ``transfer'' only happens once, i.e.\ at the time when the new environment is created, and is unidirectional, i.e.\ from the agent paired with the parent environment (\emph{parent agent}) to the \emph{child environment} with the hope that the parent agent helps jump-start learning in the child environment. 

%\ks{We should avoid talking about transfer as being from an agent to an agent.  A transfer in our system is always from an agent to an environment, completely replacing what was there before in the target environment.}\rw{agree}

In contrast, POET maintains parallel paths of environment evolution and conducts transfer experiments periodically that give every active agent multiple chances to transfer its learned skills to others, including the transfer of a child agent back to its parent environment, which continually creates and preserves opportunities to harness unanticipated stepping stones. Figure \ref{fig:transfer} presents an interesting example from POET where such transfers help a parent environment acquire a better solution: The parent environment includes only flat ground and the parent agent learns to move forward without fully standing up (top left), which works in this simple environment, but represents a local optimum in walking behavior because standing would be better. At iteration 400, the parent environment mutates and generates a child environment with small stumps, and the child agent inherits the walking gait from the parent agent. 

\begin{figure}
  \centering
     \begin{subfigure}[b]{0.9\textwidth}
        \includegraphics[width=\textwidth]{figures/transfer_story_0106}
        \caption{Transfer from agent in parent environment to child environment and vice versa}
        \label{fig:transfer_2stage}
    \end{subfigure}
    ~ %add desired spacing between images, e. g. ~, \quad, \qquad, \hfill etc. 
      %(or a blank line to force the subfigure onto a new line)
    \begin{subfigure}[b]{0.45\textwidth}
        \includegraphics[width=\textwidth]{figures/transfer_last_result_0106}
        \caption{The walking gait of agent in parent environment at Iteration 2,300}
        \label{fig:transfer_before_after}
    \end{subfigure}
\caption{\textbf{Synergistic two-way transfer between parent and child environments.} At iteration 400, a transfer from parent environment yields a child agent now learning in a stumpy environment, shown in the top row of (a). The child agent eventually learns to stand up and jump over the stumps, and at iteration 1,175 that skill is transferred \emph{back} to the parent environment, depicted in the bottom row of (a). This transfer from the child environment back to the parent helps the parent agent learn a more optimal walking gait compared to its original one. Given the same amount of computation, the agent with the original walking gaits reaches a score of 309, while the one with the more optimal walking gait as illustrated in (b) reaches a score of 349.
%\ks{The images chosen in the graphic could be better.  First, in the top row, the child at iteration 400 looks like it is *already* standing up, doesn't it?  But it is supposed to have a problem because it's croucing.  It maybe be more understandable to show it in a position of croucching next to the stump (implying it's stuck).  Second, the parent at iteration 1175 look's like it still crouching near its knee.  Shouldn't it look instead more like the parent shown at the bottom row right, fully standing?  That would be more clear and less ambiguous (it does not need to be identical to the bottom row-right image, but should be similarly upright).\rw{done. images updated.}
%\jc{ this figure is confusing. at first I figured that the left column was the flat environment and the right column was the Stumpy environment but that does not seem to hold up across each row. we need more label to which environment is which the arrows indicate.  also does not clear if we were talking about parent and children environments or parent and children agents.  also why isn't he transfer asterisk in the first row question mark isn't there a transfer background environment to the stump ground environment?  also can we add the scores of these agents to see Improvement}
%\ks{We can think about whether it can be improved.}
%\rw{Done. as we discussed, both figure and caption updated.}
%\rw{we can add scores. we can add asterisk to (a) top.for (b) we can just show parent at 2300, since we already have parent (before) in (a).  That will make it less confusing. Shall we?}
%\ks{The graphic could be improved.  Shouldn't the left side of the middle row show the agent *standing*?  The fact it is sitting makes it looks like the transfer failed to keep the character of the child.  The third row is inconsistent with the other two, both of which show transfers.  Maybe the third row should say "before" and "after" below the respective depictions to make it clear what these are, and that these are not transfers like in the other two rows.}\rw{done}
}
  \label{fig:transfer}
\end{figure}


  
 
   





At first, it can move forward in the new stumpy environment, but it often stumbles because of its crouching posture (top right). Later, through several hundred further iterations of ES, the child agent eventually learns to stand up and jump over the stumps.
The interesting moment is iteration 1,175, when that skill is transferred back to the \emph{parent} environment (middle row).  For the first time, the parent environment (which is completely flat) now contains an agent who stands up while walking.   Without the backward transfer, the agent in the simple flat environment would be stuck at its original gait (which is in effect a local optimum) forever, which was confirmed by separately turning off transfer and running the kneeling (local optimum) agent in the original (parent) environment for 3,000 iterations, which did not result in any substantive change.
%\jc{ how do we now? Did we test this hypothesis? we should  and include the results here}\ks{Rui: What kind of evidence have we seen that the kneeling gait would stay the same forever?}. \rw{for that run, i turn off transfer for that niche and observed it is on kneeling gait up to 3000 iterations.} 
Interestingly, after transfer back to the parent environment, the transferred agent then continued to optimize the standing gait to the parent (stump-free) environment and eventually obtained a much better walking gait that moves faster and costs less energy due to less friction (bottom row). More specifically, given 3,000 iterations, the agent with the original kneeling walking gait
%\jc{which original walking gait? the kneeling gait? or the standing gait just after transfer back to the parent env? in either case, how could this have the same amount of computation given that the post-transfer-and-optimize gait has the extra post-transfer optimization steps?} 
reaches a score of 309, while the one with the more optimal walking gait illustrated in figure \ref{fig:transfer_before_after} reaches a score of 349.
%\jc{how much faster?}
%\ks{Rui: any way to say something on how much faster the agent is after the transfer and subsequent optimization than before?}\rw{as we discussed, score reported here at in figure caption.}
%\jc{ so we have just pointed out one example of this happening to the reader but we leave them with no discussion of how frequent this sort of event is.  what we should do and I would as a reader expect as some sort of an analysis of how often such transfers from a parent to a child and back to a parent occur.  then more generally I would expect to hear about how often transfers  in general happen irrespective of whether they are from parent to child. in general the paper feels light on rigorous, sound quantitative analyses.}

%\ks{(new paragraph here)transfer analysis in in the works.  Rui: how is that coming along?}\rw{propose to add:} 
%\sout{\ks{Old version}
%To give a more holistic view of the success of transfer throughout the entire system, we count the number of \emph{replacement attempts} during the course of a run, which means the number of times an environment took a group of incoming transfer attempts from all the other active environments.
%The total number of such replacement attempts in RUN 1, RUN 2, and RUN 3 (labelled in Figure \ref{fig:control_type_1}) are 18,894, 19,014, and 18,798, respectively, out of which, $53.62\%$, $49.26\%$, $48.89\%$, respectively, are successful replacements. 
%Note that each replacement attempt here encompasses both the direct transfer and proposal transfer attempts from all the other environments.
%\ks{Rui:sounds like a good plan. Also then I want to add a but more text to draw a conclusion from these stats:} 
%These statistics show how pervasively transfer permeates (and often adds value to) the parallel paths explored by POET.  Every successful replacement represents a case where the performance in an environment jumped beyond what the current paired agent could achieve.\todo{Mentioned in Jeff's email: need to revise this since we are not sure if transfers really add value every time. even better would be to launch the control}}

To give a holistic view of the success of transfer throughout the entire system, we count the number of \emph{replacement attempts} during the course of a run, which means the number of times an environment took a group of incoming transfer attempts from all the other active environments.
The total number of such replacement attempts in RUN 1, RUN 2, and RUN 3 (labelled in Figure \ref{fig:control_type_1}) are 18,894, 19,014, and 18,798, respectively, out of which, $53.62\%$, $49.26\%$, $48.89\%$, respectively, are successful replacements (meaning one of the tested agents is better than the current niche champion). 
Note that each replacement attempt here encompasses both the direct transfer and proposal transfer attempts. 
%\sout{from all the other environments}\jc{without removing what I have struck out, then all attempts in an iteration (e.g. 100) should count as 1, not as 100}.
%\ks{Rui:sounds like a good plan. Also then I want to add a but more text to draw a conclusion from these stats:} 
These statistics show how pervasively transfer permeates (and often adds value to) the parallel paths explored by POET.  

\begin{comment}
\begin{figure}
  \centering
  \includegraphics[width=.85\linewidth]{figures/coverage_v2}
\caption{\textbf{The coverage of environments created and solved by POET is significantly better than those by the controls runs of POET without transfer ($p < 2.2e-16$).} Larger is worse.}
  \label{fig:coverage-boxplot}
\end{figure}
\end{comment}

%\sout{
%\rw{add: To further demonstrate the overall importance of enabling transfer in the POET algorithm, we relaunched the three POET runs but with all the transfers disabled (which is referred to later as the control runs of \emph{POET without transfer}). To do so, we skip Line 13-19 in Algorithm \ref{alg:poet}, and replace Line 16 in Algorithm \ref{alg:mutate_niches} by initializing the agent paired with a child environment the same as the agent paired with its parent environment. We can then calculate the \emph{coverage} of the environments that are created and solved by POET and by the control, respectively, following a similar metric defined in \cite{lehman:gecco11}. We first uniformly sample 1,000 \emph{challenging}, 1,000 \emph{very challenging}, and 1000 \emph{extremely challenging} environments. For each of the totally 3,000 sampled environments, the distance to the nearest environment created and solved by POET (and by control, respectively) is calculated. The better covered the environment space is, the less the sum of all such nearest distances will be. It is observed that the coverage of environments created and solved by POET is significantly better than that by the controls runs of POET without transfer ($p < 2.2e-16$) (Fig. \ref{fig:coverage-boxplot}), and that in the control without transfer, no \emph{extremely challenging} environments are generated (Fig. \ref{fig:disable-transfer-barplot}).   }
%}

While transfer is pervasive, that does not in itself prove it is essential. 
%\jc{A transfer may be successful because the incoming agent got lucky and performed better than the previous occupant by chance rather than because it is truly better (here, chance is not in the environment, which is deterministic, but instead is from noise in the optimization process).}\ks{Seems a bit complicated to try to compress into one sentence.  Probably reads better without it.}\jc{without it we never explain why transfers don't indicate that transfers are helpful. since we only accept transfers if they are better, that's pretty confusing. }  
To demonstrate the value of transfer, a control is needed: We relaunched another three POET runs, but with all the transfers \emph{disabled} (which we call \emph{POET without transfer}). In this variant, POET runs as usual, but simply never
tries to transfer paired solutions from one environment to another.
%To do so, lines 16-22 in Algorithm \ref{alg:poet} (in SI) are skipped, and line 16 in Algorithm \ref{alg:mutate_niches} (in SI) is replaced by always initializing the agent paired with a child environment as a clone of the agent paired with the parent environment.\jc{very, very confusing, and relies on SI (which ken does not like). Can we more intuitively explain this and refer to SI for extra details if the reader wants them?} \ks{agreed}
We can then calculate the \emph{coverage} of the environments that are created and solved by POET and the control, respectively, following a similar metric as defined in \citet{lehman:gecco11}: We first uniformly sample 1,000 \emph{challenging}, 1,000 \emph{very challenging}, and 1,000 \emph{extremely challenging} environments. For each of the total 3,000 sampled environments, the distance to the nearest
one of the \emph{challenging}, \emph{very challenging}, or \emph{extremely challenging} environments created and solved by POET (and by the control, respectively) is calculated. Note that the better covered the environment space is, the lower the sum of all such nearest distances will be. The result is that the coverage of environments created and solved by POET is significantly greater than by POET without transfer ($p < 2.2e{-16}$ based on Mann-Whitney U test), meaning POET with transfer explores more of the space of environments, including creating (and solving) more difficult environments. 
%(Fig. \ref{fig:coverage-boxplot}).
%(Fig. \ref{fig:coverage-boxplot}).
In an even more dramatic statistic showing the essential role of transfer in open-ended search, in POET without transfer, \emph{no extremely challenging environments are solved at all (Fig. \ref{fig:disable-transfer-barplot}).
%\jc{generated or }\ks{but it might have generated one that was not solved, depending what you mean by ``generatted''}solved at all}.\jc{exactly. the current writing is ambiguous. I was hoping it is true that it neither generated nor solves them. If true (and Rui should let us know), this correction is stronger and resolves the ambiguity. But if it is not accurate, then we should be clearer by letting the reader know it generated them, but did not solve them. I'll bet it's just that it never generated them. Note: by generated I mean generated and accepted...bleh. that's a complex distinction that I do not know if we introduce} 
%Figure \ref{fig:disable-transfer-barplot} in SI Section \ref{si:transfer-ex} illustrates this stark contrast.   
}

\begin{figure}
  \centering
  \includegraphics[width=.85\linewidth]{figures/poet_transfer_dist}
\caption{\textbf{Percentage of environments of different difficult levels created and solved by POET and the control runs of POET without transfer.} In the control without transfer, no \emph{extremely challenging} environments are generated.
}
  \label{fig:disable-transfer-barplot}
\end{figure}

%\ks{We might not want to keep this example (or its accompanying figure).  It seems like an incomplete account.  If it took transfer from a "more stumpy" environment to get the agent in this environment to go over high stumps, then how did the agents in those more stumpy ones learn to do it?  They didn't get it from a transfer did they? So why did this one need to?  And then there's the issue of the gaps.  How could a transfer from a "more stumpy" environment also know how to jump over pits (as seem to exist in this environment) unless they also could get over pits already?  But if they could, then they must have been able to get over both stumps and pits, in which case they already have solved the problem of the present environment, which makes it seem superfluous and just kicks the question back to those other environments -- how then did they learn it?  Something seems to be missing in this story.}\rw{agree. commented out the figure and context}

\begin{comment}
Figure \ref{fig:transfer2} illustrates a similar example. Transfers from less stumpy environments help the agent quickly navigate to the middle of a course, but the agent is blocked by a stump that was higher than those seen in the parent environments (left panel). However, different agents in more stumpy environments subsequently learn to jump over higher stumps, and transfers from those environments help the agent in the depicted environment gain the ability to jump higher and eventually conquer the offending stump (right panel).


\begin{figure}
  \centering
  \includegraphics[width=.75\linewidth]{figures/transfer2.png}
\caption{\textbf{Competing transfers.} Left: Transfers from less stumpy environments help the agent move smoothly over pits in this heterogeneous environment until being blocked by a higher stump that was not seen in the parent environments. Right: Transfers from a more stumpy environment supply the skill of jumping higher and eventually conquering the blocking stump. \ks{in figure 3 we use (a) and (b) to label panels; why here is it left and right?  we should be consistent.}}
  \label{fig:transfer2}
\end{figure}
\end{comment}

%\ks{To be safe, please make sure we define "solved" clearly somewhere above.  It's used constantly in this section.}\rw{Thanks for the remind. the definition is added to section 4.1}

\begin{figure}
  \centering
  \includegraphics[width=.8\linewidth]{figures/diversity_envs_2}
\caption{\textbf{Sub-sections of environments from the middle column (RUN 2) of figure \ref{fig:control_type_1}.} These environments exhibit diversity in the levels and ranges of different obstacles. For example, the topmost environment has all large gaps and high stumps with a high amount of roughness; the second from the bottom exhibits a wide variety in gap width, but with trivial stumps and low roughness. In contract, the environment at the very bottom exhibits a wide range of stump heights, but with low gap width and low roughness.
}
  \label{fig:diverse_envs}
\end{figure}

Finally, one of the primary hypothesized benefits of POET is its ability to produce a broad diversity of different problems with functional solutions in a single run. The environments (depicted as red pentagons) shown in each column of Figure \ref{fig:control_type_1} are created and solved in a \emph{single} run of POET. 
Each such column exhibits  diversity in the values and/or value ranges in roughness, gap width of gaps, and height of stumps. Take the environments in the middle column (labelled ``RUN 2'') for example. In figure \ref{fig:diverse_envs},
each image of a section of an environment from top to bottom corresponds to the same row of rose plot for RUN 2 in Figure \ref{fig:control_type_1}.
For instance, the topmost environment has narrow ranges and high values in gap width and stump height, and high value in roughness; in stark contrast the plot at the bottom depicts a less challenging environment with a wide range of stump heights, but low gap width and roughness.  
%\jc{what do these environments look like if I see pictures of them?  can we add them to supplement information?  and or make videos of them? } 
This diversity of environments also implies diverse experiences for the agents paired with them, who in turn thereby learn diverse walking gaits.  This diversity then supplies the stepping stones that fuel the mutual transfer mechanism of POET.
For a more objective statement of POET's ability to cover much of the possibility space through its diversity, all three POET runs created and solved sufficiently diverse environments to cover all three challenge levels.

%\jc{ again this discussion just creates for me an expectation that I will then see an analysis of what is happening within these runs. how  many different types of environments are generated?  how diverse are they? how does complexity change over time?  we could do an analysis of the median number of environments generated per run, and/or also in supplementary information show The Rose plot for every environment generated in a single run (e.g. run 2),  as well as show a pic of every environment generated in a single run in SI.  otherwise we could take one run I'm too and then really show that user a lot about what happens within that one run all in supplementary information.  Etc I think the reader still has little idea how much real diversity there is within a run. we're too dependent on this one plot and the few examples that we included per run.} \ks{Rui: Fair point by Jeff that it would be interesting to see more of the diversity somehow.  I don't think a ton of Rose plots is very effective because they are quite hard to interpret.  But maybe we can show screenshots of a bunch of example environments?  Also worth thinking about any kind of summary statement we could make about diversity, like the range of combinations it covers in one run (or each run).  We should discuss.}\rw{I agree that tons of Rose plots are not effective at all. Let's discuss on }
%\ks{This last paragraph is important but currently confusing because it does not refer to the diversity within a column.  Given that a column is the most important grouping in figure 4 (because a column comes from a single run), the discussion of diversity is most important within each column.  That is, we want to claim that each run is diverse.  Diversity across runs is less interesting.  The paragraph however talks about "row" instead of column, even though the meaning of a row here is vague at best (given that a row contains environments from separate independent runs).  Should be reworded to be more clear and persuasive.}\rw{agree. reworded.}
%\ks{I don't necessarily mean to imply rose plots are not effective at all. I think the ones in figure 4 do a fine job.  I just don't think we'll get that much more understanding by seeing more of them.  I also don't think we can really get all that much more visually than a few shots (of anything, rose, screenshot, whatever) because in reality this environment is limited and doesn't offer all that much diversity.  So in that sense I think we don't have to do too much to show some images.  However, what would be cool is some kind of aggregate statistic, but that is tricky to think of.}\rw{agree that the environment limit its diversity.}\ks{in any case, I do agree that we should add more screenshots of environments to qualitatively show their diversity.  the paper is missing images of almost any environments with more than one obstacle type. I think figure 3b is the only one.} \rw{if we add those screen-shots of environments, I guess it should be format like figure 1. long and narrow? it takes time to do it. How many we need? I guess we can try to plot 6 of them according to RUN 2 in Fig. 5? And I do not know how to just show the environment. The visualizer of OpenAI Gym will have the agent in there as well, i guess that is fine. So do we commit to this as todo?}\ks{Yes I think that's a good plan to give a sampling of imagery.  The 6 environments of run 2 in (now) figure 5 is a good choice.  I don't see any problem with having the agent in the shot either. We may put such a figure in SI or in the results section, but that's easy to switch.  The hard part is to make the figure.  The point of the images is to show their typical ``look'' so they shots have to be taken carefully so they aren't just some trivial moment with flat ground or something like that.} \rw{OK, let us do that. we need to make some decision about the time. i plan to give up the bar plots to show transfer statistics and just report the percentage of how much transfer are succeeded out of all transfer experiments, because bar plot basically show that percentage number. Comments?} \rw{also please comment on new Figure 4}

%\ks{what is a "representative environment"?  how were these chosen?  how many environments were there total and what percentage of the total are these?  are they chosen by hand?  are they from across the whole run or only the end?  what justifies showing them as opposed to any others?  what are they "representative" of?}\ks{Yes, more plots is good, but also once again we need a sense of how many runs there were total so we can understand the extent to which what we're seeing is anecdotal or comprehensive.}\rw{text and plots are modified to address those comments}

%\begin{figure}
%  \centering
%  \includegraphics[width=.75\linewidth]{figures/diversity_run_plot.png}
%\caption{\textbf{Representative environments that POET created and solved in one run.} Each %environment is represented as a pentagon in the plot with the same convention as in Figure %\ref{fig:control_type_1}. The diversity is evident in the different shapes of the various %environment pentagon plots.}
%  \label{fig:diversity}
%\end{figure}

%\vspace{-0.3cm}
\section{Discussion, Future Work, and Conclusion}
%\vspace{-0.23cm}

The promise of open-ended computation is fascinating for its potential to enable systems that become more powerful and interesting the longer they run.  There is always the possibility that the feats we observe in the system today will be overshadowed by the achievements of tomorrow. 
%\jl{The today/tomorrow bit is a bit confusing -- is it a literal claim that if you run POET longer it will learn to descend sequences of body-sized chasms -- I assume if we had a video of that we'd show it, or that there's a technical reason that's impossible currently, i.e. the encoding doesn't support it? I believe you intend something more metaphorical, but because it is so close to what the system already does, I'm not sure it will land.} \jc{ I also find this confusing}
In these unfolding odysseys there is an echo of the natural world.  More than just the story of a single intelligent lifetime, they evoke the history of invention, or of natural evolution over the eons of Earth.  These are processes that produce not just a single positive result, but an ongoing cacophony of surprises -- unplanned advances rolling ahead in parallel without any final destination.  

POET is an attempt to move further down the road towards these kinds of systems.  While the road remains long, the rewards for machine learning of beginning to capture the character of open-ended processes is potentially high.  First, as the results show, there is the opportunity to discover capabilities that could not be learned in any other way, even through a carefully crafted curriculum targeted at the desired result.  In addition, a diversity of such results can be generated in a single run, and the problems and solutions can both increase in complexity over time.  Furthermore, the implicit result is that POET is self-generating multiple curricula simultaneously, all while leveraging the results of some as stepping stones to progress in others.   

The experiment in this article in 2-D walking establishes the potential of POET, but POET becomes more interesting the more unbounded its problem space becomes.  The present problem space is a 2-D course of obstacles that can be generated within distributions defined by the genome describing the environment.  This space is limited by the maximal
ranges of those distributions.  For example, there is a maximum possible gap width and stump height, which means that the system in effect can eventually ``max out'' the difficulty.  While sufficiently broad to demonstrate POET's functionality, in the future much more flexible or even unbounded encodings can enable POET to traverse a far richer problem space.  For example, an indirect encoding like a compositional pattern-producing network (CPPN) \cite{stanley:gpem07} can generate arbitrarily complex patterns that can be e.g.\ converted into levels or obstacle courses.  Such an encoding would allow POET to diverge across a much richer landscape of possibilities.

An additional constraint that exists in this initial work is that the body of the agent is fixed, ultimately limiting the sort of obstacles it can overcome (e.g.\ how big of a gap it can jump over). Here too a more powerful and expressive encoding of morphologies, including those like CPPNs that are based on developmental biology, could allow us to co-evolve the morphology and body of the agent along with its brain in addition to the environment it is solving. Previous research has shown that using such indirect encodings to evolve morphologies can improve performance and create a wide diversity of interesting solutions \cite{cheney:gecco13}. Even within the 2-D Bipedal Walker domain from this paper, research has shown that allowing the morphology to be optimized can improve performance and provide a diversity of interesting solutions \cite{davidha:arxiv18}.

Furthermore, the  CPPN indirect encoding can also encode the neural network controllers in an algorithm called HyperNEAT \cite{stanley:alife09}.  Work evolving robot controllers with HyperNEAT has shown it can produce neural networks with regular geometric patterns in neural connectivity that exploit regularities in the environment to produce regular behaviors that look natural and improve performance \cite{clune:ieeetec11, clune:cec09, clune:cec09, yosinski:ecal11}. POET allows us to create the interesting combination of using an indirect encoding like a CPPN to encode the environment, the body of the agent, and the neural network controller of that agent, allowing each of the three to benefit from and exploit the regularities being produced by the others.

While ES is the base optimization algorithm in this paper for training solvers for various tasks under POET, including the main transfer mechanism, POET can be instantiated with any RL or optimization algorithm.   Trust Region Policy Optimization (TRPO) \cite{schulman2015trust}, Proximal Policy Optimization (PPO) \cite{ppo:arxiv17}, genetic algorithms, other variants of ES, and many other such algorithms are all viable alternatives.
%\jc{ Wireless just these two and not genetic algorithms, CMA-ES, etc?}. 
This plug-and-play aspect of the overall framework presents an intriguing set of future opportunities to hybridize the divergent search across environments with the most appropriate inner-loop learning algorithms, opening up the investigation of open-endedness itself to diverse research communities.

POET could also substantially drive progress in the field of meta-learning, wherein neural networks are exposed to many different problems and get better over time at learning how to solve new challenges (i.e.\ they learn to learn). Meta-learning requires access to a distribution of different tasks, and that traditionally requires a human to specify this task distribution, which is costly and may not be the right or best distribution on which to learn to learn. Gupta et al. \cite{gupta2018unsupervised} note both that the performance of meta-learning algorithms critically depends on the distribution of tasks they meta-train on, and that automatically producing an effective distribution of tasks for meta-learning algorithms would represent a substantial advance. That work took an important initial step in that direction by automatically creating different reward functions within one fixed environment \cite{gupta2018unsupervised}. POET takes the important further step of creating entirely new environments (or challenges, more abstractly). While the version of POET in this initial work does not also also create a unique reward function for each environment, doing so is a natural extension that we leave for future work. One possibility is that the reward function could be encoded in addition to the rest of the environment and optimized similarly in parallel.  

Another promising avenue for exploration in POET is its interaction with generalization versus specialization.  Interestingly, the environments in POET themselves can foster generalization by including multiple challenges in the same sequence.  However, alternate versions of POET are also conceivable that take a more explicit approach to optimizing for generality; for example, compound environments could be created whose scores are the aggregate of scores in their constituent environments (borrowing a further aspect of the CMOEA algorithm \cite{huizinga2016does, Joost:arxiv18}).  

Finally, it is exciting to consider for the future the rich potential for surprise in all the possible domains where POET might be applied.  For example, 3-D parkour was explored by \citet{heess:arxiv17} in environments created by humans, but POET could invent its own creative parkour challenges and their solutions.  The soft robots evolved by \citet{cheney:gecco13} would also be fascinating to combine with ever-unfolding new obstacle courses.  POET also offers practical opportunities in domains like autonomous driving, where through generating increasingly challenging and diverse scenarios it could uncover important edge cases and policies to solve them.  Perhaps more exotic opportunities could be conceived, such as protein folding for specific biological challenges invented by the system, or searching for new kinds of chemical processes that solve unique problems.  The scope is broad for imagination and creativity in the application of POET.  

\section*{Acknowledgments}

We thank all of the members of Uber AI Labs, in particular Joost Huizinga and Zoubin Ghahramani for helpful discussions.
We also thank Alex Gajewski (from his internship at Uber AI Labs) for help with implementation and useful discussions.   
Finally, we thank Justin
Pinkul, Mike Deats, Cody Yancey, Joel Snow, Leon Rosenshein and the entire OpusStack Team inside Uber for providing our computing platform and for technical support.


%\jc{ I feel like we should also mention that we're excited to apply this too much harder more interesting domains with a rich potential to surprise us and then list a lot of examples, especially the parkour examples from deepmind where we can say that there it was humans that created the challenges but it would be interesting to see what would happen if poet created the challenges and the solutions.  we could also say that it would be interesting to see what the soft robots of Cheney et al could do when applied to arbitrary obstacle courses created by poet.  what sorts of interesting, unique challenges could such soft robots overcome?  et cetera. This would also be a good opportunity to plug the practical uses such as in self-driving ( creating Unique Edge cases and policies that can solve them) and other things might potentially be interesting like maybe protein folding to create proteins that solves unique biological challenges, or  to create chemical processes that solve unique challenges like producing certain chemicals or  producing energy and novel ways, what else?}

%\jc{ we also haven't mentioned anything about generalists and Specialists and how we ultimately are interested in neural networks aren't one-trick ponies but can solve a wide distribution of tasks and that we think that's possible with a more powerful environment encoding with the current algorithm. however, another version of the algorithm could do something cmoea like and combine Atomic environments in a set that need to be solved which is another interesting possibility for future research. see the first todo for more motivation and things to say on this front}

\begin{comment}

OLD VERSION

In this paper, we propose POET, a novel open-ended coevolution-based framework. POET is capable of coevolving tasks with their solvers, so that increasingly difficult tasks along with increasingly capable solvers can emerge in a progressive fashion. Minimal criteria controls the evolution that create new tasks proper for the capabilities of solvers, while solvers are optimized with a novel multi-objective ES, where solvers periodically transfers their learned capability to peers to help others' learning progress. An open-endedness framework like POET provides RL researchers a new way to design/evolve a curriculum as a series of tasks with increasing complexity and difficultly, and to create, preserve, and utilize stepping stones in the optimization process toward solving progressively difficult RL problems.

We demonstrate the superior capability of POET through the example of creating and evolving 2D walking courses and the corresponding bipedal walkers. Note that this 2D environment has a direct encoding with limited number of genomes, fixed mutation rates, as well as physical limits posed by the simulation (such as predefined maximum torques on the joint), all of which cap the total number of different environments it can have. However, POET is a general framework that is not tied to any encoding and constraints above. For example, on the task side, with a richer encoding for tasks, such as an indirect encoding through compositional pattern-producing networks \cite{stanley:gpem07}, POET will likely generate even more interesting and diverse walking courses and solvers. On the solver side, though not explored in this paper, POET is fully compatible with RL algorithms that learn to modify the parameters of the agent’s body structure along with its policy \cite{davidha:arxiv18}.

Finally, in this paper we have only explored using ES as the base optimization algorithm for training solvers for various tasks under POET, and also described our transfer mechanism in terms of ES. In principle, POET can be extended to be compatible with the use of other state-of-the-art model-free RL algorithms, such as Trust Region Policy Optimization (TRPO) \cite{schulman2015trust} and Proximal Policy Optimization (PPO) \cite{ppo:arxiv17} algorithms for policy learning for agents. Then it would be interesting to explore differences in performance and behaviours in agents learning with various different RL algorithms and the impact of enabling the transfer of learned capabilities between those agents.
\end{comment}

\begin{comment}
\section{To Do}
\begin{itemize}
    \item Did we mention the other possible way of doing it? CMOEA style with niches being all possible combinations of atoms? I.e. the original way? One reason to mention it is that, while it is true (and interesting) that single environments can contain all types of environments with the right encoding, it still may be less likely. In the spirit of ``optimize for what you want your algorithm to produce'', we may be better off using these interesting atoms as a training ground for the production of generalists. \ks{added to discussion}
    \item should we mention the billion-year framing? I love that \ks{added to the intro}
    \item should we mention using ES for the environments as future work? probably not? \ks{yeah I don't think it is important for us to mention it here}
    \item Not for the first paper, but for later versions: we should show phylogenetic trees of the environments. Also add directional arcs for each transfer. Possibly color environments by performance and make a movie of the tree growing over time and the transfers radiating out cool performance lifts? \ks{would certainly be cool; significant work involved}
\end{itemize}
\end{comment}

%\subsubsection*{Acknowledgments}
%We thank all of the members of Uber AI Labs, in particular Thomas Miconi, Rui Wang, Peter Dayan, John Sears, Joost Huizinga, and Theofanis Karaletsos, for helpful discussions. We also thank Justin Pinkul, Mike Deats, Cody Yancey, Joel Snow, Leon Rosenshein and the entire OpusStack Team inside Uber for providing our computing platform and for technical support.
\bibliographystyle{unsrtnat}{\small\bibliography{nips,ucf,nn,nnstrings}}

\newpage
\section{Supplemental Information}\label{supplemental}

The first section in this supplement details the procedure in POET for generating and choosing new environments, which is followed by an explanation of how POET attempts to transfers agents from one environment to another.

%\vspace{-0.25cm}
\subsection{Producing New Environments}\label{mutation_section}
%\vspace{-0.1cm}



%\ks{Should we give an example?  The ideas here so far are a bit abstract.  It’s possible we could help the reader by saying, “For example,” and using our upcoming domain to explain the idea more concretely.  That is, niches could be obstacles courses and solvers could be neural networks that control a robot to solve the courses.  If we say something of that nature, we might both make the algorithm more concrete and also create anticipation (and eventually fulfillment) of the experiment we are soon to introduce.  The downside is the added text and growth in length of an already-long algorithmic exposition.}\rw{moved to the beginning of section, example given}

%\jc{this writing is not great. we set up in the first sentence or two the expectation that we will be discussion mutation only...but the list jumps into a much more complex process covering the whole of when and how new environments produce offspring environments, whether those are kept, what children get paired with them, etc.}
New environments are produced by copying the current environments and making random changes to them 
(in the language of evolutionary algorithms, the environments \emph{reproduce} through \emph{mutation}). 
Algorithm \ref{alg:mutate_niches}
formalizes the single mutation step, given a list of environment-agent pairs. The primary logic in this process is: (1) Each environment evaluates its paired agent. When the reward obtained by the agent satisfies a predefined reproduction eligibility condition for the environment (denoted as ELIGIBLE\_TO\_REPRODUCE()), the environment is added to the list of parent environments.
%\jc{when is this list acted on? immediately or in batches?} 
(2) Once the list of eligible parents is complete, mutations of the parent environments return a list of child environments (through the function ENV\_REPRODUCE()). Each child environment's paired agent is initialized as a clone of its parent environment's paired agent.
%\ks{this sentence from this point onward doesn’t make sense.  What does “case-by-case” mean here?  *What* is “depending on” task and encoding?  This language is very confusing because “task” is defined as a member of the list, but you’re suddenly using it differently.  Also, how did encoding suddenly come up?  The reader hasn’t heard anything about it yet so it’s totally foreign at this point.} \ks{The next sentence is a digression that seems to throw us off the trail of the explanation here.  Why should the reader suddenly care what the encoding is in the middle of explain this subroutine?}\rw{rewritten} \ks{if you have a superscript or subscript greater than 1 letter long, you should use backslash textrm to alter the font to roman from italics for that term, e.g. “child” here}\rw{corrected} 
(3) Each child environment $E^{\textrm{child}}(\cdot)$, together with its paired agent $\theta^{\textrm{child}}$, is checked against a predefined minimal criterion (reminiscent of the minimal criterion in the MCC algorithm \cite{brant:mcc17}) and only those satisfying the minimal criterion are kept (denoted as MC\_SATISFIED()), ensuring that potential new environments are not too trivial and not too hard (which would lead to little progress). 
(4) All remaining child environments are then ranked based on their novelty (detailed shortly), %\ks{need to explain in next (new) paragraph why and how novelty is calculated; completely unclear from this sentence what its purpose is or how it works.}\rw{added below.}),
which is calculated with respect to the current active population of environments and an \emph{archive} of all previously active environments (i.e.\ those that were allowed to enter the population previously).
%\jc{here is another case where archive is a more natural word. also, we should be clear: are we only comparing to environments that were at some point added to the population, or to all generated environments, including those rejected because they did not survive the MCC filter?}.
(5) A transfer attempt is then performed to see whether an agent from a different environment could lead to the best performance in a new child environment.  The logic of transfer is described in the next section (Section \ref{Transfer Overview}).  Each pair is then checked against the minimal criterion again in case some pairs no longer satisfy the minimal criterion after the transfer (e.g.\ if 
an environment is exposed as too easy). 
%POET pairs each child environment with a $\theta^{\textrm{child}}$ that is the output of Algorithm \ref{alg:es_transfer}, which takes all the current agents in the population as input candidates. The pair is then checked against \jc{the} minimal criterion again in case some pairs no longer satisfy the minimal criterion after the transfer.\jc{why do this again?} 

To limit the usage of computational resources, (6) POET caps the number of child environments that can be created during one iteration of POET's main loop.
%\jc{unclear how there could be multiple...seems like it is one at a time?} at $(\texttt{max\_children})$, 
The number of new environments that can be admitted is $\texttt{max\_admitted}$ and the maximum number of active environments at any time is $\texttt{capacity}$. If adding new environments exceeds the system's capacity, the oldest environments and their paired agents in the system are removed to make room (also as in the MCC algorithm).
%\jc{in general this whole paragraph was very rough and could be written much more clearly. 6 is the hardest to follow.}

The \emph{encoding} of an environment in POET is its underlying description in the system, such as a vector of numbers that specify the distribution of obstacles in an obstacle course.
As mentioned in step (4) above, POET calculates the novelty of the child environments.  This novelty calculation is based on the encoding. 
%\jc{we have slipped into using the language of evolution without explaining why...and that is alien to the RL/ML/DL community. also, we are using complex terms for the first time as if we have already talked about them, like genetic encoding. for each term, we should define it/introduce the reader to it}and ranks these 
Child environments are ranked based on their novelty.
For readers familiar with the \emph{novelty search} algorithm \cite{lehman:ecj11}, it is important to note here that the encoding of an environment is in effect a genetic representation, which means that in POET in effect novelty is a \emph{genetic} diversity measure.  That way, POET does not in principle require a behavior characterization (BC) \cite{lehman:ecj11} to be devised for environments, though doing so is also not precluded.
%\ks{We might want to think about language here -- "novelty" is usually associated with behavioral diversity and not genetic diversity, but in this case because the genetic encoding is a direct description of the environment, it effectively *is* the "behavior". We may want to explain this issue.  For now let's just leave this comment so we remember.}\rw{agree.}\jc{agree we should quickly explain the issue}

Measuring novelty introduces an explicit pressure for diversity among the environments by encouraging newly admitted environments to have notably different encodings than those previously admitted.
%\jc{assuming we explain the issue above that ultimately what we care about is behavioral/functional diversity, then here we should drop `genetic' and just mention that we are creating a pressure for environments to be meaningfully/behaviorally/functionally different} 
Given an environment $E$ with an encoding denoted as $e(E)$, and a list $L$
%\jc{first use of the word archive} 
containing the encodings of all the environments currently in the population plus those previously admitted but no longer active (which are stored in an \emph{archive}), 
%but removed later\jc{first mention that we remove things from an archive...no explanation of why. I think what you mean is that we have an active population and an archive... I also think you are to define archive here,  but that should be his own sentence where we tell the reader that is what's going on  instead of doing it as a parenthetical in the middle of a different complicated explanation}, 
the novelty of $E$, denoted as $N(E, L)$, is then computed by selecting the k-nearest neighbors of $e(E)$ from $L$ and computing the average distance between them:
%\jc{ does it make sense that the function n can either take the encoding of an environment or the environment itself?}
%\vspace{-0.1cm}
\begin{gather*}
\scalebox{1.}{$N(e(E), L) = \frac{1}{\vert{S}\vert}\sum_{j\in{S}}\vert\vert{e(E) - e(E_{j})}\vert\vert_{2}$}\\
\scalebox{1.}{$S = \textrm{kNN}(e(E), L)$}\\
\scalebox{1.}{$
= \{e(E_{1}), e(E_{2}), ..., e(E_{k})\}$}
\end{gather*}
Note that in this work the distance between genetic encodings of environments is calculated with an $L2$-norm and the number of nearest neighbors $k$ in the calculation of novelty is 5.

%\ks{make sure that the italicized format for variables in the algorithm pseudocode in this paper follows convention.  Also check capitalization: why are some variables capitalized and some not?}\rw{corrected}

\begin{algorithm}[htp]
\caption{MUTATE\_ENVS}
\label{alg:mutate_niches}
%\algsetup{linenosize=\tiny}
%\scriptsize
%\ks{I don't think MC\_SATISFIED is ever defined or explained anywhere}\rw{corrected}
\begin{algorithmic}[1]
   \STATE {\bfseries Input:} list of environment(E)-agent(A) pairs $\texttt{EA\_list}$,
   %\jc{confusing name because evolutionary algorithms are often called EA}, 
   each entry is a pair $(E(\cdot), \theta)$
   \STATE {\bfseries Parameters:} maximum number of children per reproduction \texttt{max\_children}, maximum number of children admitted per reproduction \texttt{max\_admitted}, maximum number of active environments \texttt{capacity}
   \STATE {\bfseries Initialize:} Set \texttt{parent\_list} to empty
   \STATE $M = \textrm{len}(\texttt{EA\_list})$   
   \FOR{$m=1$ {\bfseries to} $M$}
   \STATE $E^{m}(\cdot)$, $\theta^{m}$ = $\texttt{EA\_list}$[m]
   \IF{ELIGIBLE\_TO\_REPRODUCE$(E^{m}(\cdot), \theta^{m})$}
   \STATE add $(E^{m}(\cdot), \theta^{m})$ to \texttt{parent\_list}
   \ENDIF   
   \ENDFOR 
   \STATE $\texttt{child\_list} = $ ENV\_REPRODUCE$(\texttt{parent\_list}, \texttt{max\_children})$%\jc{where do we find out how maxChildren works in this function? which envs are preferred if the cap limits things?}
   \STATE $\texttt{child\_list} = $ MC\_SATISFIED$(\texttt{child\_list})$%\jc{did we define this function?}
   \STATE $\texttt{child\_list} = $ RANK\_BY\_NOVELTY$(\texttt{child\_list})$
   \STATE $\texttt{admitted} = 0$
   \FOR{$E^{\textrm{child}}(\cdot), \theta^{\textrm{child}} \in \texttt{child\_list}$}
   \STATE $\theta^{\textrm{child}} = $ EVALUATE\_CANDIDATES$(\theta^{1}$, $\theta^{2}$, \ldots, $\theta^{M}$, $E^{\textrm{child}}(\cdot)$, $\alpha$, $\sigma))$
   \IF{MC\_SATISFIED$([E^{\textrm{child}}(\cdot), \theta^{\textrm{child}}])$} 
   \STATE add $(E^{\textrm{child}}(\cdot), \theta^{\textrm{child}})$ to $\texttt{EA\_list}$
   \STATE $\texttt{admitted} = \texttt{admitted} + 1$
   \STATE {\bfseries if} $\texttt{admitted} >= \texttt{max\_admitted}$ {\bfseries then} break {\bfseries end if}%\jc{what is the effect of this ordering? which type of envs tend to get excluded?}
   \ENDIF
   \ENDFOR
   \STATE $M = \textrm{len}(\texttt{EA\_list})$
   \IF{$M > \texttt{capacity}$}
   \STATE $\texttt{num\_removals} = M - \texttt{capacity}$
   \STATE REMOVE\_OLDEST$(\texttt{EA\_list}, \texttt{num\_removals})$ 
   \ENDIF
   \STATE {\bfseries Return:} $\texttt{EA\_list}$


   
\end{algorithmic}
\end{algorithm}

\subsection{ES-Based Transfer} \label{Transfer Overview}
%\vspace{-0.2cm}
POET adopts a transfer mechanism that allows one agent (the \emph{emigrant})
%\jc{emmigrant?} 
to attempt to transfer its learned capabilities to 
%\ks{I removed text that introduced the concept of \emph{target agent} because that term misleadingly suggests that information from one agent is being transfered into another agents, which is not the case.  Rather, the agent in the target environment is completely replaced if the transfer works out.  New language aims to be clearer in this issue, but note we've lost the term \emph{target agent} from our vocabulary so we need to make sure it isn't used in the future.}\rw{agree.}
%another agent (the \emph{target agent}) to boost the learning progress in the 
a \emph{target environment}, possibly replacing the current target environment's agent. This transfer mechanism reflects the idea that skills gained when interacting with one environment might serve as a stepping stone to improved %\jc{improved} 
performance when interacting with another environment.  We do not know at any given time which environment's mastery might boost another's, but we can always perform \emph{experiments} to see whether such a transfer might improve performance in the target environment. As results support, the richness of simultaneous transfer attempts in a divergent coevolving system can lead to a succession of breakthroughs that would not be possible through a simple monotonic succession of improvements on a single environment.

%\ks{The paragraph below is unclear and confusing: I don’t understand the first sentence.  The “origin solver” is just one parameter vector (theta) right?  So how can the “transferred parameter vector” be “the better one” when there is only one?  It also seems to say that the “origin solver’s parameter vector” is “referred to as ‘direct transfer’”, which doesn’t make sense: how can the *vector* be called “direct transfer?”  Also, emphasis (backslash emph) is better than quotes when defining terms.  Quotes are for euphemisms, and we don’t want to imply that these technical terms are euphamisms.  Overall, paragraph needs to be rewritten.}\rw{rewritten}

%\ks{Should we really use F to denote the tasks?  A task is really more than a fitness function.  it could be called an \emph{environment} or a \emph{niche} or even a \emph{task}.  By using the variable F we seem to be saying that the fitness function is what varies, but actually, the fitness function is the same in all environments isn't it?  It's the environment that changes as opposed to the fitness function.  I'm also not sure if we want \emph{fitness} so prominent here when actually the whole system could work with something other than ES too, where the word fitness wouldn't be as likely to be used.}\rw{I change the notation from $F$ to $E$ throughout the paper}

%HERE2

Algorithm \ref{alg:es_transfer} illustrates the function EVALUATE\_CANDIDATES() that facilitates this transfer process. Given a list of $M$ input candidate agents, denoted as $\theta^{1}$, $\theta^{2}$, \ldots, $\theta^{M}$, and a target environment $E(\cdot)$, this function first concatenates the list with \emph{proposals}, which are calculated by optimizing each of $\theta^{1}$, $\theta^{2}$, \ldots, $\theta^{M}$ in $E(\cdot)$ with one ES step. Then each agent in the augmented list (the original $M$ agents and each of their $M$ proposals) is evaluated in $E(\cdot)$ and the one with the highest reward is returned. As illustrated in Algorithm \ref{alg:poet}, the agent returned from this function will replace the agent currently paired with $E(\cdot)$ if the former performs better than the latter in $E(\cdot)$. We refer to the replacement as \emph{transfer}, and there are two types: the situation when the replacing agent is from the input candidate list (i.e. without an optimization step on $E(\cdot)$) is referred to as a \emph{direct transfer}, while if the replacing agent is one of the proposals, it is called a \emph{proposal transfer}. 
%\ks{The way this paragraph ends seems to imply that whichever candidate ranks highest on F will replace whatever was there before, but actually whatever theta was already the current best for F might still be better. In other words, the current occupant of the niche might survive all the proposals.  I realize we haven’t introduced the concept of niches yet, but can we do anything at the end of this paragraph to anticipate the possibility that none of the attempted transfers ends up making it through because what’s already there is good enough to stay the best?  Or maybe the idea here is that the theta list *includes* the current theta for F, but then that one isn’t really a “transfer” right? (And the alg is called ES\_TRANSFER.)  So we could be fostering a little confusion here if the language isn’t careful.}\rw{rewritten}

\begin{algorithm}[htp]
\caption{EVALUATE\_CANDIDATES}
\label{alg:es_transfer}
%\algsetup{linenosize=\tiny}
%\scriptsize
\begin{algorithmic}[1]
   \STATE {\bfseries Input:} candidate agents denoted by their policy parameter vectors $\theta^{1}$, $\theta^{2}$, \ldots, $\theta^{M}$, target environment $E(\cdot)$,  learning rate $\alpha$, noise standard deviation $\sigma$
   \STATE {\bfseries Initialize:} Set list $C$ empty 
   \FOR{$m=1$ {\bfseries to} $M$}
   \STATE Add $\theta^{m}$ to $C$
   \STATE Add $(\theta^{m} +$ ES\_STEP$(\theta^{m}, E(\cdot), \alpha, \sigma)) $ to $C$
   \ENDFOR
   %\STATE {\bfseries Return:} $E(\theta)$
   \STATE {\bfseries Return:} $\argmax_{\theta \in \mathcal{C}} E(\theta) $

   
\end{algorithmic}
\end{algorithm}



\end{document}